% paper70-natural-language-specs.tex — Bridging Natural Language Specifications
%   to Formal Judgments
%   Paper 70 of the Judgment Geometry series.
% Compile: pdflatex paper70-natural-language-specs
\documentclass[11pt]{article}
\usepackage{lmodern}
% jugeo-common.tex — shared preamble for all Judgment Geometry papers
% Include via: % jugeo-common.tex — shared preamble for all Judgment Geometry papers
% Include via: % jugeo-common.tex — shared preamble for all Judgment Geometry papers
% Include via: \input{jugeo-common}

% ── Packages ────────────────────────────────────────────────────────────
\usepackage{amsmath,amssymb,amsthm,mathtools}
\usepackage{tikz-cd}
\usepackage{listings}
\usepackage{xcolor}
\usepackage{xspace}
\usepackage{booktabs}
\usepackage{multirow}
\usepackage{enumitem}
\usepackage[expansion=false]{microtype}
\usepackage{hyperref}
\usepackage{cleveref}
% \usepackage{thmtools}  % disabled: not available in this TeX installation
\usepackage{float}
\usepackage{caption}
\usepackage{subcaption}
\usepackage{tabularx}
\usepackage{stmaryrd}       % \llbracket, \rrbracket
\usepackage{bbm}            % \mathbbm{1}

% ── Page geometry ───────────────────────────────────────────────────────
\usepackage[margin=1in]{geometry}

% ── Theorem environments ────────────────────────────────────────────────
\theoremstyle{plain}
\newtheorem{theorem}{Theorem}[section]
\newtheorem{lemma}[theorem]{Lemma}
\newtheorem{proposition}[theorem]{Proposition}
\newtheorem{corollary}[theorem]{Corollary}

\theoremstyle{definition}
\newtheorem{definition}[theorem]{Definition}
\newtheorem{example}[theorem]{Example}
\newtheorem{construction}[theorem]{Construction}

\theoremstyle{remark}
\newtheorem{remark}[theorem]{Remark}
\newtheorem{notation}[theorem]{Notation}

% ── Python listings style ──────────────────────────────────────────────
\definecolor{jugeoBlue}{HTML}{2563EB}
\definecolor{jugeoGreen}{HTML}{059669}
\definecolor{jugeoRed}{HTML}{DC2626}
\definecolor{jugeoPurple}{HTML}{7C3AED}
\definecolor{jugeoGray}{HTML}{6B7280}
\definecolor{jugeoBg}{HTML}{F8FAFC}

\lstdefinestyle{jugeo-python}{
  language=Python,
  basicstyle=\ttfamily\small,
  keywordstyle=\color{jugeoBlue}\bfseries,
  stringstyle=\color{jugeoGreen},
  commentstyle=\color{jugeoGray}\itshape,
  numberstyle=\tiny\color{jugeoGray},
  backgroundcolor=\color{jugeoBg},
  numbers=left,
  numbersep=6pt,
  frame=single,
  framerule=0.4pt,
  rulecolor=\color{jugeoGray!40},
  breaklines=true,
  breakatwhitespace=true,
  showstringspaces=false,
  tabsize=4,
  xleftmargin=1.5em,
  framexleftmargin=1.5em,
  morekeywords={dataclass,frozen,field,Enum,IntEnum,auto,Any,
                Mapping,Sequence,Optional,match,case},
  literate={→}{{$\to$}}1 {←}{{$\leftarrow$}}1
           {≤}{{$\leq$}}1 {≥}{{$\geq$}}1 {≠}{{$\neq$}}1
           {∈}{{$\in$}}1 {∀}{{$\forall$}}1 {∃}{{$\exists$}}1
           {⊕}{{$\oplus$}}1 {⊖}{{$\ominus$}}1,
  escapeinside={(*@}{@*)},
}
\lstset{style=jugeo-python}

\lstdefinestyle{jugeo-output}{
  basicstyle=\ttfamily\small,
  backgroundcolor=\color{jugeoBg},
  frame=single,
  framerule=0.4pt,
  rulecolor=\color{jugeoGray!40},
  breaklines=true,
  showstringspaces=false,
  xleftmargin=1.5em,
  framexleftmargin=0.5em,
  numbers=none,
}

% ── JuGeo-specific macros ──────────────────────────────────────────────

% Judgment 8-tuple
\newcommand{\J}{\mathcal{J}}
\newcommand{\Jud}[8]{(#1,\,#2,\,#3,\,#4,\,#5,\,#6,\,#7,\,#8)}
\newcommand{\JudShort}[1]{\J_{#1}}

% Components
\newcommand{\coord}{\mathsf{c}}            % coordinate
\newcommand{\prop}{\varphi}                 % proposition
\newcommand{\carrier}{\mathcal{A}}          % carrier type
\newcommand{\evid}{\mathcal{E}}             % evidence bundle
\newcommand{\oblig}{\mathcal{O}}            % obligations
\newcommand{\obstr}{\mathcal{B}}            % obstructions
\newcommand{\trust}{\mathcal{T}}            % trust annotation
\newcommand{\prov}{\Pi}                     % provenance

% Trust algebra
\newcommand{\Talg}{\mathfrak{T}}            % trust ordered algebra
\newcommand{\Eadm}{\mathcal{E}_{\mathrm{adm}}}
\newcommand{\tleq}{\preceq}                 % trust ordering
\newcommand{\tjoin}{\oplus}                 % conservative join
\newcommand{\tatten}{\ominus}               % attenuation
\newcommand{\tpromote}{\uparrow_\pi}        % promotion
\newcommand{\tdemote}{\downarrow_\chi}       % demotion

% Trust tiers
\newcommand{\Tcontra}{\ensuremath{\mathsf{CONTRADICTED}}}
\newcommand{\Tunver}{\ensuremath{\mathsf{UNVERIFIED}}}
\newcommand{\Tcopilot}{\ensuremath{\mathsf{COPILOT}}}
\newcommand{\Toracle}{\ensuremath{\mathsf{ORACLE}}}
\newcommand{\Truntime}{\ensuremath{\mathsf{RUNTIME}}}
\newcommand{\Tsolver}{\ensuremath{\mathsf{SOLVER}}}
\newcommand{\Tproof}{\ensuremath{\mathsf{PROOF}}}

% Site and geometry
\newcommand{\Site}{\mathcal{S}}             % semantic site
\newcommand{\Cov}{\mathrm{Cov}}            % covering family
\newcommand{\Ob}{\mathrm{Ob}}              % objects
\newcommand{\Mor}{\mathrm{Mor}}            % morphisms
\newcommand{\res}{\mathrm{res}}            % restriction
\newcommand{\Cech}{\check{C}}              % Čech complex
\newcommand{\Hcoh}[1]{H^{#1}}             % cohomology group

% Descent
\newcommand{\descent}{\mathrm{desc}}
\newcommand{\glue}{\mathrm{glue}}
\newcommand{\overlap}[2]{#1 \cap #2}

% Semantic moves
\newcommand{\VERIFY}{\textsc{Verify}}
\newcommand{\CONSTRUCT}{\textsc{Construct}}
\newcommand{\NORMALIZE}{\textsc{Normalize}}
\newcommand{\GLUE}{\textsc{Glue}}
\newcommand{\RESTRICT}{\textsc{Restrict}}
\newcommand{\TRANSPORT}{\textsc{Transport}}
\newcommand{\REPAIR}{\textsc{Repair}}
\newcommand{\PROMOTE}{\textsc{Promote}}
\newcommand{\CHALLENGE}{\textsc{Challenge}}
\newcommand{\ESCALATE}{\textsc{Escalate}}

% Fragments
\newcommand{\QFLIA}{\texttt{QF\_LIA}}
\newcommand{\QFLRA}{\texttt{QF\_LRA}}
\newcommand{\QFBV}{\texttt{QF\_BV}}
\newcommand{\QFUF}{\texttt{QF\_UF}}

% Misc
\newcommand{\Python}{\textsc{Python}}
\newcommand{\JuGeo}{\textsc{JuGeo}}
\newcommand{\LEAN}{\textsc{Lean}\,4}
\newcommand{\Fstar}{F$^{\star}$}
\newcommand{\Coq}{\textsc{Coq}}
\newcommand{\Dafny}{\textsc{Dafny}}

% Common shortcuts
\newcommand{\ie}{i.e.\@\xspace}
\newcommand{\eg}{e.g.\@\xspace}
\newcommand{\cf}{cf.\@\xspace}
\newcommand{\wrt}{w.r.t.\@\xspace}

% Evaluation shortcuts
\newcommand{\numcases}{300}
\newcommand{\numspec}{100}
\newcommand{\numeq}{100}
\newcommand{\numbug}{100}
\newcommand{\medtime}{6.5\,ms}
\newcommand{\totaltime}{1.949\,s}


% ── Packages ────────────────────────────────────────────────────────────
\usepackage{amsmath,amssymb,amsthm,mathtools}
\usepackage{tikz-cd}
\usepackage{listings}
\usepackage{xcolor}
\usepackage{xspace}
\usepackage{booktabs}
\usepackage{multirow}
\usepackage{enumitem}
\usepackage[expansion=false]{microtype}
\usepackage{hyperref}
\usepackage{cleveref}
% \usepackage{thmtools}  % disabled: not available in this TeX installation
\usepackage{float}
\usepackage{caption}
\usepackage{subcaption}
\usepackage{tabularx}
\usepackage{stmaryrd}       % \llbracket, \rrbracket
\usepackage{bbm}            % \mathbbm{1}

% ── Page geometry ───────────────────────────────────────────────────────
\usepackage[margin=1in]{geometry}

% ── Theorem environments ────────────────────────────────────────────────
\theoremstyle{plain}
\newtheorem{theorem}{Theorem}[section]
\newtheorem{lemma}[theorem]{Lemma}
\newtheorem{proposition}[theorem]{Proposition}
\newtheorem{corollary}[theorem]{Corollary}

\theoremstyle{definition}
\newtheorem{definition}[theorem]{Definition}
\newtheorem{example}[theorem]{Example}
\newtheorem{construction}[theorem]{Construction}

\theoremstyle{remark}
\newtheorem{remark}[theorem]{Remark}
\newtheorem{notation}[theorem]{Notation}

% ── Python listings style ──────────────────────────────────────────────
\definecolor{jugeoBlue}{HTML}{2563EB}
\definecolor{jugeoGreen}{HTML}{059669}
\definecolor{jugeoRed}{HTML}{DC2626}
\definecolor{jugeoPurple}{HTML}{7C3AED}
\definecolor{jugeoGray}{HTML}{6B7280}
\definecolor{jugeoBg}{HTML}{F8FAFC}

\lstdefinestyle{jugeo-python}{
  language=Python,
  basicstyle=\ttfamily\small,
  keywordstyle=\color{jugeoBlue}\bfseries,
  stringstyle=\color{jugeoGreen},
  commentstyle=\color{jugeoGray}\itshape,
  numberstyle=\tiny\color{jugeoGray},
  backgroundcolor=\color{jugeoBg},
  numbers=left,
  numbersep=6pt,
  frame=single,
  framerule=0.4pt,
  rulecolor=\color{jugeoGray!40},
  breaklines=true,
  breakatwhitespace=true,
  showstringspaces=false,
  tabsize=4,
  xleftmargin=1.5em,
  framexleftmargin=1.5em,
  morekeywords={dataclass,frozen,field,Enum,IntEnum,auto,Any,
                Mapping,Sequence,Optional,match,case},
  literate={→}{{$\to$}}1 {←}{{$\leftarrow$}}1
           {≤}{{$\leq$}}1 {≥}{{$\geq$}}1 {≠}{{$\neq$}}1
           {∈}{{$\in$}}1 {∀}{{$\forall$}}1 {∃}{{$\exists$}}1
           {⊕}{{$\oplus$}}1 {⊖}{{$\ominus$}}1,
  escapeinside={(*@}{@*)},
}
\lstset{style=jugeo-python}

\lstdefinestyle{jugeo-output}{
  basicstyle=\ttfamily\small,
  backgroundcolor=\color{jugeoBg},
  frame=single,
  framerule=0.4pt,
  rulecolor=\color{jugeoGray!40},
  breaklines=true,
  showstringspaces=false,
  xleftmargin=1.5em,
  framexleftmargin=0.5em,
  numbers=none,
}

% ── JuGeo-specific macros ──────────────────────────────────────────────

% Judgment 8-tuple
\newcommand{\J}{\mathcal{J}}
\newcommand{\Jud}[8]{(#1,\,#2,\,#3,\,#4,\,#5,\,#6,\,#7,\,#8)}
\newcommand{\JudShort}[1]{\J_{#1}}

% Components
\newcommand{\coord}{\mathsf{c}}            % coordinate
\newcommand{\prop}{\varphi}                 % proposition
\newcommand{\carrier}{\mathcal{A}}          % carrier type
\newcommand{\evid}{\mathcal{E}}             % evidence bundle
\newcommand{\oblig}{\mathcal{O}}            % obligations
\newcommand{\obstr}{\mathcal{B}}            % obstructions
\newcommand{\trust}{\mathcal{T}}            % trust annotation
\newcommand{\prov}{\Pi}                     % provenance

% Trust algebra
\newcommand{\Talg}{\mathfrak{T}}            % trust ordered algebra
\newcommand{\Eadm}{\mathcal{E}_{\mathrm{adm}}}
\newcommand{\tleq}{\preceq}                 % trust ordering
\newcommand{\tjoin}{\oplus}                 % conservative join
\newcommand{\tatten}{\ominus}               % attenuation
\newcommand{\tpromote}{\uparrow_\pi}        % promotion
\newcommand{\tdemote}{\downarrow_\chi}       % demotion

% Trust tiers
\newcommand{\Tcontra}{\ensuremath{\mathsf{CONTRADICTED}}}
\newcommand{\Tunver}{\ensuremath{\mathsf{UNVERIFIED}}}
\newcommand{\Tcopilot}{\ensuremath{\mathsf{COPILOT}}}
\newcommand{\Toracle}{\ensuremath{\mathsf{ORACLE}}}
\newcommand{\Truntime}{\ensuremath{\mathsf{RUNTIME}}}
\newcommand{\Tsolver}{\ensuremath{\mathsf{SOLVER}}}
\newcommand{\Tproof}{\ensuremath{\mathsf{PROOF}}}

% Site and geometry
\newcommand{\Site}{\mathcal{S}}             % semantic site
\newcommand{\Cov}{\mathrm{Cov}}            % covering family
\newcommand{\Ob}{\mathrm{Ob}}              % objects
\newcommand{\Mor}{\mathrm{Mor}}            % morphisms
\newcommand{\res}{\mathrm{res}}            % restriction
\newcommand{\Cech}{\check{C}}              % Čech complex
\newcommand{\Hcoh}[1]{H^{#1}}             % cohomology group

% Descent
\newcommand{\descent}{\mathrm{desc}}
\newcommand{\glue}{\mathrm{glue}}
\newcommand{\overlap}[2]{#1 \cap #2}

% Semantic moves
\newcommand{\VERIFY}{\textsc{Verify}}
\newcommand{\CONSTRUCT}{\textsc{Construct}}
\newcommand{\NORMALIZE}{\textsc{Normalize}}
\newcommand{\GLUE}{\textsc{Glue}}
\newcommand{\RESTRICT}{\textsc{Restrict}}
\newcommand{\TRANSPORT}{\textsc{Transport}}
\newcommand{\REPAIR}{\textsc{Repair}}
\newcommand{\PROMOTE}{\textsc{Promote}}
\newcommand{\CHALLENGE}{\textsc{Challenge}}
\newcommand{\ESCALATE}{\textsc{Escalate}}

% Fragments
\newcommand{\QFLIA}{\texttt{QF\_LIA}}
\newcommand{\QFLRA}{\texttt{QF\_LRA}}
\newcommand{\QFBV}{\texttt{QF\_BV}}
\newcommand{\QFUF}{\texttt{QF\_UF}}

% Misc
\newcommand{\Python}{\textsc{Python}}
\newcommand{\JuGeo}{\textsc{JuGeo}}
\newcommand{\LEAN}{\textsc{Lean}\,4}
\newcommand{\Fstar}{F$^{\star}$}
\newcommand{\Coq}{\textsc{Coq}}
\newcommand{\Dafny}{\textsc{Dafny}}

% Common shortcuts
\newcommand{\ie}{i.e.\@\xspace}
\newcommand{\eg}{e.g.\@\xspace}
\newcommand{\cf}{cf.\@\xspace}
\newcommand{\wrt}{w.r.t.\@\xspace}

% Evaluation shortcuts
\newcommand{\numcases}{300}
\newcommand{\numspec}{100}
\newcommand{\numeq}{100}
\newcommand{\numbug}{100}
\newcommand{\medtime}{6.5\,ms}
\newcommand{\totaltime}{1.949\,s}


% ── Packages ────────────────────────────────────────────────────────────
\usepackage{amsmath,amssymb,amsthm,mathtools}
\usepackage{tikz-cd}
\usepackage{listings}
\usepackage{xcolor}
\usepackage{xspace}
\usepackage{booktabs}
\usepackage{multirow}
\usepackage{enumitem}
\usepackage[expansion=false]{microtype}
\usepackage{hyperref}
\usepackage{cleveref}
% \usepackage{thmtools}  % disabled: not available in this TeX installation
\usepackage{float}
\usepackage{caption}
\usepackage{subcaption}
\usepackage{tabularx}
\usepackage{stmaryrd}       % \llbracket, \rrbracket
\usepackage{bbm}            % \mathbbm{1}

% ── Page geometry ───────────────────────────────────────────────────────
\usepackage[margin=1in]{geometry}

% ── Theorem environments ────────────────────────────────────────────────
\theoremstyle{plain}
\newtheorem{theorem}{Theorem}[section]
\newtheorem{lemma}[theorem]{Lemma}
\newtheorem{proposition}[theorem]{Proposition}
\newtheorem{corollary}[theorem]{Corollary}

\theoremstyle{definition}
\newtheorem{definition}[theorem]{Definition}
\newtheorem{example}[theorem]{Example}
\newtheorem{construction}[theorem]{Construction}

\theoremstyle{remark}
\newtheorem{remark}[theorem]{Remark}
\newtheorem{notation}[theorem]{Notation}

% ── Python listings style ──────────────────────────────────────────────
\definecolor{jugeoBlue}{HTML}{2563EB}
\definecolor{jugeoGreen}{HTML}{059669}
\definecolor{jugeoRed}{HTML}{DC2626}
\definecolor{jugeoPurple}{HTML}{7C3AED}
\definecolor{jugeoGray}{HTML}{6B7280}
\definecolor{jugeoBg}{HTML}{F8FAFC}

\lstdefinestyle{jugeo-python}{
  language=Python,
  basicstyle=\ttfamily\small,
  keywordstyle=\color{jugeoBlue}\bfseries,
  stringstyle=\color{jugeoGreen},
  commentstyle=\color{jugeoGray}\itshape,
  numberstyle=\tiny\color{jugeoGray},
  backgroundcolor=\color{jugeoBg},
  numbers=left,
  numbersep=6pt,
  frame=single,
  framerule=0.4pt,
  rulecolor=\color{jugeoGray!40},
  breaklines=true,
  breakatwhitespace=true,
  showstringspaces=false,
  tabsize=4,
  xleftmargin=1.5em,
  framexleftmargin=1.5em,
  morekeywords={dataclass,frozen,field,Enum,IntEnum,auto,Any,
                Mapping,Sequence,Optional,match,case},
  literate={→}{{$\to$}}1 {←}{{$\leftarrow$}}1
           {≤}{{$\leq$}}1 {≥}{{$\geq$}}1 {≠}{{$\neq$}}1
           {∈}{{$\in$}}1 {∀}{{$\forall$}}1 {∃}{{$\exists$}}1
           {⊕}{{$\oplus$}}1 {⊖}{{$\ominus$}}1,
  escapeinside={(*@}{@*)},
}
\lstset{style=jugeo-python}

\lstdefinestyle{jugeo-output}{
  basicstyle=\ttfamily\small,
  backgroundcolor=\color{jugeoBg},
  frame=single,
  framerule=0.4pt,
  rulecolor=\color{jugeoGray!40},
  breaklines=true,
  showstringspaces=false,
  xleftmargin=1.5em,
  framexleftmargin=0.5em,
  numbers=none,
}

% ── JuGeo-specific macros ──────────────────────────────────────────────

% Judgment 8-tuple
\newcommand{\J}{\mathcal{J}}
\newcommand{\Jud}[8]{(#1,\,#2,\,#3,\,#4,\,#5,\,#6,\,#7,\,#8)}
\newcommand{\JudShort}[1]{\J_{#1}}

% Components
\newcommand{\coord}{\mathsf{c}}            % coordinate
\newcommand{\prop}{\varphi}                 % proposition
\newcommand{\carrier}{\mathcal{A}}          % carrier type
\newcommand{\evid}{\mathcal{E}}             % evidence bundle
\newcommand{\oblig}{\mathcal{O}}            % obligations
\newcommand{\obstr}{\mathcal{B}}            % obstructions
\newcommand{\trust}{\mathcal{T}}            % trust annotation
\newcommand{\prov}{\Pi}                     % provenance

% Trust algebra
\newcommand{\Talg}{\mathfrak{T}}            % trust ordered algebra
\newcommand{\Eadm}{\mathcal{E}_{\mathrm{adm}}}
\newcommand{\tleq}{\preceq}                 % trust ordering
\newcommand{\tjoin}{\oplus}                 % conservative join
\newcommand{\tatten}{\ominus}               % attenuation
\newcommand{\tpromote}{\uparrow_\pi}        % promotion
\newcommand{\tdemote}{\downarrow_\chi}       % demotion

% Trust tiers
\newcommand{\Tcontra}{\ensuremath{\mathsf{CONTRADICTED}}}
\newcommand{\Tunver}{\ensuremath{\mathsf{UNVERIFIED}}}
\newcommand{\Tcopilot}{\ensuremath{\mathsf{COPILOT}}}
\newcommand{\Toracle}{\ensuremath{\mathsf{ORACLE}}}
\newcommand{\Truntime}{\ensuremath{\mathsf{RUNTIME}}}
\newcommand{\Tsolver}{\ensuremath{\mathsf{SOLVER}}}
\newcommand{\Tproof}{\ensuremath{\mathsf{PROOF}}}

% Site and geometry
\newcommand{\Site}{\mathcal{S}}             % semantic site
\newcommand{\Cov}{\mathrm{Cov}}            % covering family
\newcommand{\Ob}{\mathrm{Ob}}              % objects
\newcommand{\Mor}{\mathrm{Mor}}            % morphisms
\newcommand{\res}{\mathrm{res}}            % restriction
\newcommand{\Cech}{\check{C}}              % Čech complex
\newcommand{\Hcoh}[1]{H^{#1}}             % cohomology group

% Descent
\newcommand{\descent}{\mathrm{desc}}
\newcommand{\glue}{\mathrm{glue}}
\newcommand{\overlap}[2]{#1 \cap #2}

% Semantic moves
\newcommand{\VERIFY}{\textsc{Verify}}
\newcommand{\CONSTRUCT}{\textsc{Construct}}
\newcommand{\NORMALIZE}{\textsc{Normalize}}
\newcommand{\GLUE}{\textsc{Glue}}
\newcommand{\RESTRICT}{\textsc{Restrict}}
\newcommand{\TRANSPORT}{\textsc{Transport}}
\newcommand{\REPAIR}{\textsc{Repair}}
\newcommand{\PROMOTE}{\textsc{Promote}}
\newcommand{\CHALLENGE}{\textsc{Challenge}}
\newcommand{\ESCALATE}{\textsc{Escalate}}

% Fragments
\newcommand{\QFLIA}{\texttt{QF\_LIA}}
\newcommand{\QFLRA}{\texttt{QF\_LRA}}
\newcommand{\QFBV}{\texttt{QF\_BV}}
\newcommand{\QFUF}{\texttt{QF\_UF}}

% Misc
\newcommand{\Python}{\textsc{Python}}
\newcommand{\JuGeo}{\textsc{JuGeo}}
\newcommand{\LEAN}{\textsc{Lean}\,4}
\newcommand{\Fstar}{F$^{\star}$}
\newcommand{\Coq}{\textsc{Coq}}
\newcommand{\Dafny}{\textsc{Dafny}}

% Common shortcuts
\newcommand{\ie}{i.e.\@\xspace}
\newcommand{\eg}{e.g.\@\xspace}
\newcommand{\cf}{cf.\@\xspace}
\newcommand{\wrt}{w.r.t.\@\xspace}

% Evaluation shortcuts
\newcommand{\numcases}{300}
\newcommand{\numspec}{100}
\newcommand{\numeq}{100}
\newcommand{\numbug}{100}
\newcommand{\medtime}{6.5\,ms}
\newcommand{\totaltime}{1.949\,s}

% experiment-data.tex — AUTO-GENERATED from real jugeo CLI runs
% DO NOT EDIT — regenerate with: python3 experiments/run_universal_metrics.py
% Generated from 30 programs across 6 domains

\newcommand{\expTotalPrograms}{30}
\newcommand{\expVerified}{30}
\newcommand{\expAccuracy}{100.0\%}
\newcommand{\expCoordsMin}{2}
\newcommand{\expCoordsMax}{10}
\newcommand{\expCoordsMean}{5.2}
\newcommand{\expCoordsMedian}{5.5}
\newcommand{\expPropsMin}{4}
\newcommand{\expPropsMax}{20}
\newcommand{\expPropsMean}{10.4}
\newcommand{\expPropsSum}{311}
\newcommand{\expPropsOkSum}{310}
\newcommand{\expTimeMin}{3.506\,s}
\newcommand{\expTimeMax}{6.714\,s}
\newcommand{\expTimeMean}{4.64\,s}
\newcommand{\expTimeMedian}{4.508\,s}
\newcommand{\expTimeTotal}{139.204\,s}
\newcommand{\expObstructionTotal}{0}
\newcommand{\expHOneZero}{30}
\newcommand{\expDomainCount}{6}
\newcommand{\expTrustCOPILOTSUGGESTED}{30}
\newcommand{\expDomainDsCount}{5}
\newcommand{\expDomainDsVerified}{5}
\newcommand{\expDomainDsAvgCoords}{7.6}
\newcommand{\expDomainDsAvgProps}{15.2}
\newcommand{\expDomainMathCount}{5}
\newcommand{\expDomainMathVerified}{5}
\newcommand{\expDomainMathAvgCoords}{4.8}
\newcommand{\expDomainMathAvgProps}{9.4}
\newcommand{\expDomainSmCount}{5}
\newcommand{\expDomainSmVerified}{5}
\newcommand{\expDomainSmAvgCoords}{7.6}
\newcommand{\expDomainSmAvgProps}{15.2}
\newcommand{\expDomainSortCount}{5}
\newcommand{\expDomainSortVerified}{5}
\newcommand{\expDomainSortAvgCoords}{2.4}
\newcommand{\expDomainSortAvgProps}{4.8}
\newcommand{\expDomainStrCount}{5}
\newcommand{\expDomainStrVerified}{5}
\newcommand{\expDomainStrAvgCoords}{3.2}
\newcommand{\expDomainStrAvgProps}{6.4}
\newcommand{\expDomainWebCount}{5}
\newcommand{\expDomainWebVerified}{5}
\newcommand{\expDomainWebAvgCoords}{5.6}
\newcommand{\expDomainWebAvgProps}{11.2}

% subsystem-data.tex — AUTO-GENERATED from real jugeo subsystem experiments
% DO NOT EDIT — regenerate with: python3 experiments/run_subsystem_experiments.py

\newcommand{\subCliTotal}{10}
\newcommand{\subCliVerified}{9}
\newcommand{\subCliAccuracy}{90.0\%}
\newcommand{\subCliMeanTime}{4.132\,s}
\newcommand{\subCliMinTime}{3.472\,s}
\newcommand{\subCliMaxTime}{5.372\,s}
\newcommand{\subCliTotalCoords}{54}
\newcommand{\subCliTotalProps}{107}
\newcommand{\subCliTotalPropsOk}{106}
\newcommand{\subSiteCalls}{90}
\newcommand{\subSiteSuccessful}{69}
\newcommand{\subSiteSuccessRate}{76.7\%}
\newcommand{\subSiteMeanTime}{0.0001\,s}
\newcommand{\subSiteTotalTime}{0.005\,s}
\newcommand{\subSpecificationsatisfactionSmallTime}{0.0003\,s}
\newcommand{\subSpecificationsatisfactionMediumTime}{0.0001\,s}
\newcommand{\subSpecificationsatisfactionLargeTime}{0.0001\,s}
\newcommand{\subInhabitantfleetSmallTime}{0.0\,s}
\newcommand{\subInhabitantfleetMediumTime}{0.0\,s}
\newcommand{\subInhabitantfleetLargeTime}{0.0\,s}
\newcommand{\subTheoremecologySmallTime}{0.0\,s}
\newcommand{\subTheoremecologyMediumTime}{0.0\,s}
\newcommand{\subTheoremecologyLargeTime}{0.0\,s}
\newcommand{\subStatespaceexplorationSmallTime}{0.0\,s}
\newcommand{\subStatespaceexplorationMediumTime}{0.0\,s}
\newcommand{\subStatespaceexplorationLargeTime}{0.0\,s}
\newcommand{\subRepairsemanticsSmallTime}{0.0\,s}
\newcommand{\subRepairsemanticsMediumTime}{0.0\,s}
\newcommand{\subRepairsemanticsLargeTime}{0.0\,s}
\newcommand{\subReplaygluingSmallTime}{0.0\,s}
\newcommand{\subReplaygluingMediumTime}{0.0\,s}
\newcommand{\subReplaygluingLargeTime}{0.0\,s}
\newcommand{\subSemanticfuturesSmallTime}{0.0\,s}
\newcommand{\subSemanticfuturesMediumTime}{0.0\,s}
\newcommand{\subSemanticfuturesLargeTime}{0.0\,s}
\newcommand{\subHypercovertreatySmallTime}{0.0\,s}
\newcommand{\subHypercovertreatyMediumTime}{0.0\,s}
\newcommand{\subHypercovertreatyLargeTime}{0.0\,s}
\newcommand{\subEncodeforsolverSmallTime}{0.0026\,s}
\newcommand{\subEncodeforsolverMediumTime}{0.0002\,s}
\newcommand{\subEncodeforsolverLargeTime}{0.0001\,s}
\newcommand{\subInterfaceroutingSmallTime}{0.0\,s}
\newcommand{\subInterfaceroutingMediumTime}{0.0\,s}
\newcommand{\subInterfaceroutingLargeTime}{0.0\,s}
\newcommand{\subOrchestrateverificationSmallTime}{0.0002\,s}
\newcommand{\subOrchestrateverificationMediumTime}{0.0\,s}
\newcommand{\subOrchestrateverificationLargeTime}{0.0\,s}
\newcommand{\subRunfulldescentSmallTime}{0.0\,s}
\newcommand{\subRunfulldescentMediumTime}{0.0\,s}
\newcommand{\subRunfulldescentLargeTime}{0.0\,s}
\newcommand{\subKernellifecycleSmallTime}{0.0\,s}
\newcommand{\subKernellifecycleMediumTime}{0.0\,s}
\newcommand{\subKernellifecycleLargeTime}{0.0\,s}
\newcommand{\subDiscoverypipelineSmallTime}{0.0\,s}
\newcommand{\subDiscoverypipelineMediumTime}{0.0\,s}
\newcommand{\subDiscoverypipelineLargeTime}{0.0\,s}
\newcommand{\subAnalogytransportSmallTime}{0.0\,s}
\newcommand{\subAnalogytransportMediumTime}{0.0\,s}
\newcommand{\subAnalogytransportLargeTime}{0.0\,s}
\newcommand{\subFormalcoresiteSmallTime}{0.0001\,s}
\newcommand{\subFormalcoresiteMediumTime}{0.0\,s}
\newcommand{\subFormalcoresiteLargeTime}{0.0\,s}
\newcommand{\subProblematlasSmallTime}{0.0\,s}
\newcommand{\subProblematlasMediumTime}{0.0\,s}
\newcommand{\subProblematlasLargeTime}{0.0\,s}
\newcommand{\subPublicalignmentSmallTime}{0.0\,s}
\newcommand{\subPublicalignmentMediumTime}{0.0\,s}
\newcommand{\subPublicalignmentLargeTime}{0.0\,s}
\newcommand{\subRegimebootstrappingSmallTime}{0.0\,s}
\newcommand{\subRegimebootstrappingMediumTime}{0.0\,s}
\newcommand{\subRegimebootstrappingLargeTime}{0.0\,s}
\newcommand{\subRelationalrefinementSmallTime}{0.0\,s}
\newcommand{\subRelationalrefinementMediumTime}{0.0\,s}
\newcommand{\subRelationalrefinementLargeTime}{0.0\,s}
\newcommand{\subSemanticclosureSmallTime}{0.0\,s}
\newcommand{\subSemanticclosureMediumTime}{0.0\,s}
\newcommand{\subSemanticclosureLargeTime}{0.0\,s}
\newcommand{\subTheoremeconomicsSmallTime}{0.0001\,s}
\newcommand{\subTheoremeconomicsMediumTime}{0.0\,s}
\newcommand{\subTheoremeconomicsLargeTime}{0.0\,s}
\newcommand{\subBenchmarksuiteSmallTime}{0.0\,s}
\newcommand{\subBenchmarksuiteMediumTime}{0.0\,s}
\newcommand{\subBenchmarksuiteLargeTime}{0.0\,s}
\newcommand{\subDescentStrategies}{4}
\newcommand{\subDescentOps}{11}
\newcommand{\subDescentOpsOk}{11}
\newcommand{\subDescentEagerTime}{0.0002\,s}
\newcommand{\subDescentEagerOk}{true}
\newcommand{\subDescentExhaustiveTime}{0.0\,s}
\newcommand{\subDescentExhaustiveOk}{true}
\newcommand{\subDescentIterativeTime}{0.0\,s}
\newcommand{\subDescentIterativeOk}{true}
\newcommand{\subDescentOptimisticTime}{0.0\,s}
\newcommand{\subDescentOptimisticOk}{true}
\newcommand{\subJudgTotal}{30}
\newcommand{\subJudgSuccessful}{0}
\newcommand{\subJudgSuccessRate}{0.0\%}
\newcommand{\subJudgMeanTimeUs}{7.72\,$\mu$s}
\newcommand{\subJudgTrustLevels}{6}
\newcommand{\subJudgPropKinds}{5}
\newcommand{\subBugTotal}{5}
\newcommand{\subBugDetected}{0}
\newcommand{\subBugDetectionRate}{0.0\%}
\newcommand{\subBugFalsePositiveRate}{0.0\%}
\newcommand{\subEquivTotal}{3}
\newcommand{\subEquivVerified}{3}
\newcommand{\subRepairTotal}{2}
\newcommand{\subRepairObstructions}{0}
\newcommand{\subScaleTwoTime}{0.0001\,s}
\newcommand{\subScaleFourTime}{0.0\,s}
\newcommand{\subScaleEightTime}{0.0\,s}
\newcommand{\subScaleTwelveTime}{0.0\,s}
\newcommand{\subScaleSixteenTime}{0.0\,s}
\newcommand{\subScaleTwentyTime}{0.0\,s}
\newcommand{\subTotalExperimentTime}{95.6\,s}

% comprehensive-data.tex — AUTO-GENERATED from real jugeo experiments
% DO NOT EDIT — regenerate with: python3 experiments/run_comprehensive_experiments.py
% Generated: 2026-03-23 09:19:06

% --- Size scaling metrics ---
\newcommand{\compTotalPrograms}{18}
\newcommand{\compTotalVerified}{18}
\newcommand{\compOverallAccuracy}{100.0\%}
\newcommand{\compTinyCount}{5}
\newcommand{\compTinyVerified}{5}
\newcommand{\compTinyMeanCoords}{3}
\newcommand{\compTinyMeanProps}{6}
\newcommand{\compTinyMeanTime}{4.95\,s}
\newcommand{\compTinyMeanLines}{5}
\newcommand{\compTinyMinTime}{4.45\,s}
\newcommand{\compTinyMaxTime}{5.51\,s}
\newcommand{\compSmallCount}{5}
\newcommand{\compSmallVerified}{5}
\newcommand{\compSmallMeanCoords}{3.6}
\newcommand{\compSmallMeanProps}{7.2}
\newcommand{\compSmallMeanTime}{4.85\,s}
\newcommand{\compSmallMeanLines}{16.0}
\newcommand{\compSmallMinTime}{4.30\,s}
\newcommand{\compSmallMaxTime}{5.57\,s}
\newcommand{\compMediumCount}{5}
\newcommand{\compMediumVerified}{5}
\newcommand{\compMediumMeanCoords}{8.6}
\newcommand{\compMediumMeanProps}{17.2}
\newcommand{\compMediumMeanTime}{4.47\,s}
\newcommand{\compMediumMeanLines}{45.0}
\newcommand{\compMediumMinTime}{4.26\,s}
\newcommand{\compMediumMaxTime}{4.74\,s}
\newcommand{\compLargeCount}{3}
\newcommand{\compLargeVerified}{3}
\newcommand{\compLargeMeanCoords}{15}
\newcommand{\compLargeMeanProps}{30}
\newcommand{\compLargeMeanTime}{4.31\,s}
\newcommand{\compLargeMeanLines}{79.0}
\newcommand{\compLargeMinTime}{4.27\,s}
\newcommand{\compLargeMaxTime}{4.38\,s}
% --- Aggregate timing ---
\newcommand{\compTimeMean}{4.68\,s}
\newcommand{\compTimeMedian}{4.50\,s}
\newcommand{\compTimeMin}{4.265\,s}
\newcommand{\compTimeMax}{5.571\,s}
\newcommand{\compTimeTotal}{84.3\,s}
\newcommand{\compTimeStdev}{0.45\,s}
% --- Aggregate structure ---
\newcommand{\compCoordsMean}{6.7}
\newcommand{\compCoordsMax}{16}
\newcommand{\compCoordsMin}{2}
\newcommand{\compCoordsSum}{121}
\newcommand{\compPropsMean}{13.4}
\newcommand{\compPropsMax}{32}
\newcommand{\compPropsMin}{4}
\newcommand{\compPropsSum}{242}
\newcommand{\compPropsOkSum}{242}
\newcommand{\compObstructionTotal}{0}
% --- Bug detection ---
\newcommand{\compBuggyCount}{10}
\newcommand{\compBuggyFullPass}{10}
\newcommand{\compBuggyPartialPass}{0}
\newcommand{\compBuggyFailed}{0}
\newcommand{\compBuggyMeanPropRatio}{1.00}
\newcommand{\compCorrectMeanPropRatio}{1.00}
\newcommand{\compBuggyMeanTime}{5.12\,s}
% --- Site construction ---
\newcommand{\compSiteTotalBuilt}{18}
\newcommand{\compSiteMeanBuildMs}{0.05}
\newcommand{\compSiteMaxBuildMs}{0.14}
\newcommand{\compSiteMeanCoords}{6.5}
\newcommand{\compSiteMeanMorphisms}{14.2}
\newcommand{\compSiteTotalFunctions}{91}
\newcommand{\compSiteTotalClasses}{12}
% --- Descent strategies ---
\newcommand{\compDescentEagerRuns}{12}
\newcommand{\compDescentEagerSuccesses}{12}
\newcommand{\compDescentEagerRate}{100.0\%}
\newcommand{\compDescentEagerMeanMs}{0.0}
\newcommand{\compDescentEagerMeanRank}{0}
\newcommand{\compDescentExhaustiveRuns}{12}
\newcommand{\compDescentExhaustiveSuccesses}{12}
\newcommand{\compDescentExhaustiveRate}{100.0\%}
\newcommand{\compDescentExhaustiveMeanMs}{0.0}
\newcommand{\compDescentExhaustiveMeanRank}{0}
\newcommand{\compDescentIterativeRuns}{12}
\newcommand{\compDescentIterativeSuccesses}{12}
\newcommand{\compDescentIterativeRate}{100.0\%}
\newcommand{\compDescentIterativeMeanMs}{0.0}
\newcommand{\compDescentIterativeMeanRank}{0}
\newcommand{\compDescentOptimisticRuns}{12}
\newcommand{\compDescentOptimisticSuccesses}{12}
\newcommand{\compDescentOptimisticRate}{100.0\%}
\newcommand{\compDescentOptimisticMeanMs}{0.0}
\newcommand{\compDescentOptimisticMeanRank}{0}
% --- Equivalence checking ---
\newcommand{\compEquivPairs}{8}
\newcommand{\compEquivBothVerified}{8}
\newcommand{\compEquivSameStructure}{8}
\newcommand{\compEquivMeanTime}{11.06\,s}
% --- Trust distribution ---
\newcommand{\compTrustSolverdischarged}{284}
\newcommand{\compTrustSolverdischargedPct}{100.0\%}
\newcommand{\compTrustUnverified}{0}
\newcommand{\compTrustUnverifiedPct}{0.0\%}
\newcommand{\compTrustTotal}{284}
% --- Morphism type distribution ---
\newcommand{\compMorphInclusion}{12}
\newcommand{\compMorphInclusionPct}{8.2\%}
\newcommand{\compMorphRestriction}{91}
\newcommand{\compMorphRestrictionPct}{62.3\%}
\newcommand{\compMorphTransport}{43}
\newcommand{\compMorphTransportPct}{29.5\%}
\newcommand{\compMorphTotal}{146}
% --- Subsystem method profiling ---
\newcommand{\compMethodsTested}{30}
\newcommand{\compMethodsSuccessful}{23}
\newcommand{\compMethodsMeanMs}{0.02}
\newcommand{\compProfAnalogyTransportMeanMs}{0.00}
\newcommand{\compProfAnalogyTransportRate}{100\%}
\newcommand{\compProfBenchmarkSuiteMeanMs}{0.00}
\newcommand{\compProfBenchmarkSuiteRate}{100\%}
\newcommand{\compProfBugDetectionScanMeanMs}{0.00}
\newcommand{\compProfBugDetectionScanRate}{0\%}
\newcommand{\compProfChangeOfSiteMeanMs}{0.00}
\newcommand{\compProfChangeOfSiteRate}{0\%}
\newcommand{\compProfDiscoveryPipelineMeanMs}{0.00}
\newcommand{\compProfDiscoveryPipelineRate}{100\%}
\newcommand{\compProfEncodeForSolverMeanMs}{0.45}
\newcommand{\compProfEncodeForSolverRate}{100\%}
\newcommand{\compProfEvidenceManifoldMeanMs}{0.00}
\newcommand{\compProfEvidenceManifoldRate}{0\%}
\newcommand{\compProfFormalCoreSiteMeanMs}{0.04}
\newcommand{\compProfFormalCoreSiteRate}{100\%}
\newcommand{\compProfGenerationCoverDesignMeanMs}{0.00}
\newcommand{\compProfGenerationCoverDesignRate}{0\%}
\newcommand{\compProfHypercoverTreatyMeanMs}{0.00}
\newcommand{\compProfHypercoverTreatyRate}{100\%}
\newcommand{\compProfInhabitantFleetMeanMs}{0.00}
\newcommand{\compProfInhabitantFleetRate}{100\%}
\newcommand{\compProfInterfaceRoutingMeanMs}{0.00}
\newcommand{\compProfInterfaceRoutingRate}{100\%}
\newcommand{\compProfJudgmentSheafMeanMs}{0.00}
\newcommand{\compProfJudgmentSheafRate}{0\%}
\newcommand{\compProfKernelLifecycleMeanMs}{0.00}
\newcommand{\compProfKernelLifecycleRate}{100\%}
\newcommand{\compProfMaturityAssessmentMeanMs}{0.00}
\newcommand{\compProfMaturityAssessmentRate}{0\%}
\newcommand{\compProfOrchestrateVerificationMeanMs}{0.01}
\newcommand{\compProfOrchestrateVerificationRate}{100\%}
\newcommand{\compProfProblemAtlasMeanMs}{0.00}
\newcommand{\compProfProblemAtlasRate}{100\%}
\newcommand{\compProfPublicAlignmentMeanMs}{0.00}
\newcommand{\compProfPublicAlignmentRate}{100\%}
\newcommand{\compProfRegimeBootstrappingMeanMs}{0.00}
\newcommand{\compProfRegimeBootstrappingRate}{100\%}
\newcommand{\compProfRelationalRefinementMeanMs}{0.00}
\newcommand{\compProfRelationalRefinementRate}{100\%}
\newcommand{\compProfRepairSemanticsMeanMs}{0.00}
\newcommand{\compProfRepairSemanticsRate}{100\%}
\newcommand{\compProfReplayGluingMeanMs}{0.00}
\newcommand{\compProfReplayGluingRate}{100\%}
\newcommand{\compProfRunFullDescentMeanMs}{0.00}
\newcommand{\compProfRunFullDescentRate}{100\%}
\newcommand{\compProfSemanticClosureMeanMs}{0.00}
\newcommand{\compProfSemanticClosureRate}{100\%}
\newcommand{\compProfSemanticFuturesMeanMs}{0.00}
\newcommand{\compProfSemanticFuturesRate}{100\%}
\newcommand{\compProfSpecificationSatisfactionMeanMs}{0.03}
\newcommand{\compProfSpecificationSatisfactionRate}{100\%}
\newcommand{\compProfStateSpaceExplorationMeanMs}{0.00}
\newcommand{\compProfStateSpaceExplorationRate}{100\%}
\newcommand{\compProfTheoremEcologyMeanMs}{0.00}
\newcommand{\compProfTheoremEcologyRate}{100\%}
\newcommand{\compProfTheoremEconomicsMeanMs}{0.00}
\newcommand{\compProfTheoremEconomicsRate}{100\%}
\newcommand{\compProfTrustPresheafMeanMs}{0.00}
\newcommand{\compProfTrustPresheafRate}{0\%}

% data-paper70.tex — AUTO-GENERATED by exp70_natural_language_specs.py
% DO NOT EDIT — regenerate with: python3 experiments/exp70_natural_language_specs.py
% Generated from 12 NL-annotated programs

\newcommand{\ppLXXprogramCount}{12}
\newcommand{\ppLXXprogramsOk}{12}
\newcommand{\ppLXXverified}{12}
\newcommand{\ppLXXverifiedPct}{100.0\%}

\newcommand{\ppLXXpropsTotal}{76}
\newcommand{\ppLXXpropsOk}{76}
\newcommand{\ppLXXpropRate}{100.0\%}
\newcommand{\ppLXXpropsMean}{6.33}

\newcommand{\ppLXXnlLinesTotal}{82}
\newcommand{\ppLXXnlLinesMean}{6.83}
\newcommand{\ppLXXcoordsMean}{3.17}

\newcommand{\ppLXXtimeMean}{20.43\,s}
\newcommand{\ppLXXtimeTotal}{245.19\,s}
\newcommand{\ppLXXtimeMin}{14.384\,s}
\newcommand{\ppLXXtimeMax}{26.137\,s}

% Per-program NL pipeline results
\newcommand{\ppLXXnlsortedlistCoords}{2}
\newcommand{\ppLXXnlsortedlistProps}{4}
\newcommand{\ppLXXnlsortedlistOk}{4}
\newcommand{\ppLXXnlsortedlistNlLines}{7}
\newcommand{\ppLXXnlsortedlistVerdict}{verified}
\newcommand{\ppLXXnlsortedlistTime}{14.384\,s}
\newcommand{\ppLXXnluniqueelementsCoords}{2}
\newcommand{\ppLXXnluniqueelementsProps}{4}
\newcommand{\ppLXXnluniqueelementsOk}{4}
\newcommand{\ppLXXnluniqueelementsNlLines}{5}
\newcommand{\ppLXXnluniqueelementsVerdict}{verified}
\newcommand{\ppLXXnluniqueelementsTime}{15.272\,s}
\newcommand{\ppLXXnlpositivefilterCoords}{2}
\newcommand{\ppLXXnlpositivefilterProps}{4}
\newcommand{\ppLXXnlpositivefilterOk}{4}
\newcommand{\ppLXXnlpositivefilterNlLines}{5}
\newcommand{\ppLXXnlpositivefilterVerdict}{verified}
\newcommand{\ppLXXnlpositivefilterTime}{22.896\,s}
\newcommand{\ppLXXnlsafedivideCoords}{2}
\newcommand{\ppLXXnlsafedivideProps}{4}
\newcommand{\ppLXXnlsafedivideOk}{4}
\newcommand{\ppLXXnlsafedivideNlLines}{5}
\newcommand{\ppLXXnlsafedivideVerdict}{verified}
\newcommand{\ppLXXnlsafedivideTime}{26.137\,s}
\newcommand{\ppLXXnlboundedbufferCoords}{6}
\newcommand{\ppLXXnlboundedbufferProps}{12}
\newcommand{\ppLXXnlboundedbufferOk}{12}
\newcommand{\ppLXXnlboundedbufferNlLines}{8}
\newcommand{\ppLXXnlboundedbufferVerdict}{verified}
\newcommand{\ppLXXnlboundedbufferTime}{25.392\,s}
\newcommand{\ppLXXnlstringvalidatorCoords}{2}
\newcommand{\ppLXXnlstringvalidatorProps}{4}
\newcommand{\ppLXXnlstringvalidatorOk}{4}
\newcommand{\ppLXXnlstringvalidatorNlLines}{6}
\newcommand{\ppLXXnlstringvalidatorVerdict}{verified}
\newcommand{\ppLXXnlstringvalidatorTime}{19.367\,s}
\newcommand{\ppLXXnlrangecheckCoords}{2}
\newcommand{\ppLXXnlrangecheckProps}{4}
\newcommand{\ppLXXnlrangecheckOk}{4}
\newcommand{\ppLXXnlrangecheckNlLines}{5}
\newcommand{\ppLXXnlrangecheckVerdict}{verified}
\newcommand{\ppLXXnlrangecheckTime}{17.548\,s}
\newcommand{\ppLXXnlfibonaccispecCoords}{2}
\newcommand{\ppLXXnlfibonaccispecProps}{4}
\newcommand{\ppLXXnlfibonaccispecOk}{4}
\newcommand{\ppLXXnlfibonaccispecNlLines}{7}
\newcommand{\ppLXXnlfibonaccispecVerdict}{verified}
\newcommand{\ppLXXnlfibonaccispecTime}{16.056\,s}
\newcommand{\ppLXXnlstackcontractCoords}{7}
\newcommand{\ppLXXnlstackcontractProps}{14}
\newcommand{\ppLXXnlstackcontractOk}{14}
\newcommand{\ppLXXnlstackcontractNlLines}{11}
\newcommand{\ppLXXnlstackcontractVerdict}{verified}
\newcommand{\ppLXXnlstackcontractTime}{16.501\,s}
\newcommand{\ppLXXnlmatrixmultiplyCoords}{2}
\newcommand{\ppLXXnlmatrixmultiplyProps}{4}
\newcommand{\ppLXXnlmatrixmultiplyOk}{4}
\newcommand{\ppLXXnlmatrixmultiplyNlLines}{5}
\newcommand{\ppLXXnlmatrixmultiplyVerdict}{verified}
\newcommand{\ppLXXnlmatrixmultiplyTime}{21.759\,s}
\newcommand{\ppLXXnlsearchwithspecCoords}{2}
\newcommand{\ppLXXnlsearchwithspecProps}{4}
\newcommand{\ppLXXnlsearchwithspecOk}{4}
\newcommand{\ppLXXnlsearchwithspecNlLines}{8}
\newcommand{\ppLXXnlsearchwithspecVerdict}{verified}
\newcommand{\ppLXXnlsearchwithspecTime}{26.062\,s}
\newcommand{\ppLXXnlcounterinvariantCoords}{7}
\newcommand{\ppLXXnlcounterinvariantProps}{14}
\newcommand{\ppLXXnlcounterinvariantOk}{14}
\newcommand{\ppLXXnlcounterinvariantNlLines}{10}
\newcommand{\ppLXXnlcounterinvariantVerdict}{verified}
\newcommand{\ppLXXnlcounterinvariantTime}{23.815\,s}


% ── Additional packages ────────────────────────────────────────────
\usepackage{array}
\usepackage{tikz}
\usetikzlibrary{arrows.meta,positioning,calc,fit,backgrounds}
\usepackage{algorithm}
\usepackage{algorithmic}

% ── Paper-specific macros ──────────────────────────────────────────
\newcommand{\NLBridge}{\textsc{NLBridge}}
\newcommand{\SpecParser}{\textsc{SpecParser}}
\newcommand{\JudgSynth}{\textsc{JudgmentSynthesizer}}
\newcommand{\OrchestrateVerif}{\texttt{orchestrate\_verification}}
\newcommand{\EncodeForSolver}{\texttt{encode\_for\_solver}}
\newcommand{\SpecSatisfaction}{\texttt{specification\_satisfaction}}
\newcommand{\TrustPresheaf}{\texttt{trust\_presheaf}}
\newcommand{\CopilotSugg}{\mathsf{COPILOT\_SUGGESTED}}
\newcommand{\SolverDisch}{\mathsf{SOLVER\_DISCHARGED}}
\newcommand{\TrustUpgrade}{\rightsquigarrow}
\newcommand{\NLSpec}{\mathcal{N}}
\newcommand{\FormalSpec}{\mathcal{F}}
\newcommand{\IntentLattice}{\mathcal{I}}
\newcommand{\SpecFunctor}{\Phi}

\title{\textbf{Bridging Natural Language Specifications\\
to Formal Judgments}}

\author{JuGeo Research Group}
\date{\today}

\begin{document}
\maketitle

% ─────────────────────────────────────────────────────────────────────
\begin{abstract}
Natural language specifications are the primary means by which
developers express requirements, yet they remain disconnected from
formal verification. We present \NLBridge\ as a design sketch for bridging
natural-language requirements to JuGeo judgments through parsing, synthesis,
and verification. The public repository audited for this paper does not expose
a stable \texttt{NLBridge} Python API, so code listings below are explicitly
treated as pseudocode for the proposed translation layer. The empirical part
of the paper is limited to the Paper~70 corpus of
\ppLXXprogramCount{} NL-annotated programs, all of which verify
(\ppLXXpropsOk{}/\ppLXXpropsTotal{} propositions) with mean runtime
\ppLXXtimeMean{}.
\end{abstract}

\newpage

% =====================================================================
\section{Introduction}\label{sec:intro}
% =====================================================================

Specifications in natural language---docstrings, comments, user
stories---contain rich semantic information inaccessible to formal
verification tools.  The gap between NL intent and formal specification
is a major barrier to practical formal methods.

\paragraph{The specification gap.}
Consider: \emph{``This function returns a sorted list with no
duplicates.''} This encodes at least three properties: permutation,
sortedness, and uniqueness.  Current tools require manual formalization.

\paragraph{The bridge functor.}
\JuGeo\ closes this gap with \NLBridge, treating translation as a
functor $\SpecFunctor : \NLSpec \to \FormalSpec$.  The source category
has NL fragments as objects and semantic entailment as morphisms; the
target has judgment sheaves and trust-preserving maps.

\paragraph{Contributions.}
\begin{itemize}[noitemsep]
  \item Intent lattice $\IntentLattice$ for structured NL
        representation (§\ref{sec:intent-lattice}).
  \item Bridge functor $\SpecFunctor$ and its properties
        (§\ref{sec:bridge-functor}).
  \item Three-stage NL-to-judgment pipeline (§\ref{sec:pipeline}).
  \item Specification Preservation Theorem (§\ref{sec:preservation}).
  \item Lean~4 proof sketches (§\ref{sec:lean-proof}).
  \item Evaluation on \ppLXXprogramCount{} NL-annotated programs
        (§\ref{sec:evaluation}).
\end{itemize}

% =====================================================================
\section{Background}\label{sec:background}
% =====================================================================

\JuGeo\ models verification as sheaf theory over $\Site$.  Each
$\coord$ carries $\J = (\coord, \varphi, P, Q, I, \delta, \tau, \pi)$.
Prior NL processing work~\cite{Harris2003,Ghosh2016} extracts
structured requirements from text.  LLM approaches~\cite{Chen2021}
generate annotations but lack compositional guarantees.  The
\SpecSatisfaction\ API checks whether judgments satisfy a specification;
combined with \OrchestrateVerif, it provides end-to-end verification.

% =====================================================================
\section{Intent Lattice}\label{sec:intent-lattice}
% =====================================================================

\begin{definition}[Intent]\label{def:intent}
An \emph{intent} $\iota \in \IntentLattice$ is a typed semantic unit
from NL.  The lattice comprises: \texttt{type\_constraint},
\texttt{range\_bound}, \texttt{ordering}, \texttt{uniqueness},
\texttt{nullability}, \texttt{side\_effect}, \texttt{resource\_bound},
\texttt{relation}, \texttt{invariant}, \texttt{exception}.
Each intent has strength: \textsf{must}, \textsf{should}, \textsf{may}.
\end{definition}

The lattice is ordered by specificity (\eg $\texttt{ordering} \leq
\texttt{ascending\_order} \leq \texttt{strictly\_ascending}$) with
join $\sqcup$ and meet $\sqcap$.

\begin{lstlisting}[style=jugeo-python]
# Illustrative pseudocode: the audited repository does not currently
# expose a stable NLBridge parsing API.
def parse_nl_spec(text: str) -> list[Intent]:
    """Parse natural language into structured intents."""
    intents = []
    patterns = {
        r"sorted|in order|ascending|descending":
            IntentType.ORDERING,
        r"no duplicates|unique|distinct":
            IntentType.UNIQUENESS,
        r"non-negative|positive|at least|at most":
            IntentType.RANGE_BOUND,
        r"not null|non-null|must exist":
            IntentType.NULLABILITY,
        r"returns? (?:a |an )?(\w+)":
            IntentType.TYPE_CONSTRAINT,
        r"invariant|always|never changes":
            IntentType.INVARIANT,
    }
    for pattern, itype in patterns.items():
        for m in re.finditer(pattern, text, re.IGNORECASE):
            strength = extract_strength(text, m.span())
            intents.append(Intent(type=itype,
                strength=strength, raw_text=m.group()))
    llm_intents = llm_extract_intents(text, intents)
    intents.extend(llm_intents)
    return deduplicate(intents)
\end{lstlisting}

\begin{proposition}[Intent Composition]\label{prop:intent-compose}
For intents $\iota_1, \iota_2$ from different sources covering
$\coord$: if $\iota_1 \leq \iota_2$, use $\iota_2$; if incomparable,
use $\iota_1 \sqcap \iota_2$; if contradictory, flag inconsistency.
\end{proposition}

% =====================================================================
\section{Bridge Functor}\label{sec:bridge-functor}
% =====================================================================

\begin{definition}[NL Specification Category]\label{def:nl-cat}
$\NLSpec$ has NL fragments as objects and semantic entailment
$n_1 \Rightarrow n_2$ as morphisms.
\end{definition}

\begin{definition}[Bridge Functor]\label{def:bridge}
$\SpecFunctor : \NLSpec \to \FormalSpec$ maps objects
$\SpecFunctor(n) = \mathcal{F}_n$ (synthesized judgment sheaf) and
morphisms $\SpecFunctor(n_1 \Rightarrow n_2)$ to the induced sheaf
morphism.
\end{definition}

\begin{theorem}[Functor Laws]\label{thm:functor-laws}
$\SpecFunctor$ preserves identity ($\SpecFunctor(\mathrm{id}_n) =
\mathrm{id}_{\SpecFunctor(n)}$) and composition
($\SpecFunctor(g \circ f) = \SpecFunctor(g) \circ \SpecFunctor(f)$).
\end{theorem}

\begin{proof}
Identity: self-entailment maps to the identity sheaf morphism.
Composition: follows from transitivity of both NL entailment and
logical implication.
\end{proof}

% =====================================================================
\section{NL-to-Judgment Pipeline}\label{sec:pipeline}
% =====================================================================

\subsection{Judgment Synthesis}

\begin{lstlisting}[style=jugeo-python]
def synthesize_judgment(coord, intents):
    """Map parsed intents to a judgment tuple."""
    pre, post, inv = [], [], []
    for intent in intents:
        if intent.type == IntentType.TYPE_CONSTRAINT:
            post.append(TypeProp(coord.return_var,
                                 intent.target_type))
        elif intent.type == IntentType.ORDERING:
            post.append(SortedProp(coord.return_var,
                                   intent.direction))
        elif intent.type == IntentType.UNIQUENESS:
            post.append(UniqueProp(coord.return_var))
        elif intent.type == IntentType.RANGE_BOUND:
            p = RangeProp(intent.variable,
                         intent.lower, intent.upper)
            (pre if intent.is_precondition else post).append(p)
        elif intent.type == IntentType.NULLABILITY:
            post.append(NonNullProp(intent.variable))
        elif intent.type == IntentType.INVARIANT:
            inv.append(InvariantProp(intent.expression))
    return Judgment(coord=coord, props=pre+post+inv,
        precond=pre, postcond=post, inv=inv,
        trust="COPILOT_SUGGESTED",
        provenance=f"NLBridge: {coord.docstring[:80]}")
\end{lstlisting}

\subsection{Verification Pipeline}

\begin{algorithm}[H]
\caption{NL-to-Verified Judgment Pipeline}
\label{alg:nl-pipeline}
\begin{algorithmic}[1]
\STATE \textbf{Input:} Program $P$ with NL specs
\STATE \textbf{Output:} Verified judgments $\mathcal{J}_v$
\STATE $\mathcal{J}_v \gets \emptyset$
\FOR{each $\coord \in \Site(P)$}
  \STATE $\text{intents} \gets \text{parse\_nl\_spec}
         (\text{extract\_nl}(\coord))$
  \STATE $\J \gets \text{synthesize\_judgment}(\coord, \text{intents})$
  \IF{$\SpecSatisfaction(\J, \coord)$}
    \STATE $\phi \gets \EncodeForSolver(\J)$
    \IF{$\text{Z3.prove}(\phi) = \text{success}$}
      \STATE $\J' \gets \J$ with $\tau = \SolverDisch$
      \STATE $\mathcal{J}_v \gets \mathcal{J}_v \cup \{\J'\}$
    \ENDIF
  \ENDIF
\ENDFOR
\RETURN $\mathcal{J}_v$
\end{algorithmic}
\end{algorithm}

Ambiguity is handled via multiple interpretations (verify all, keep
successes), confidence scoring, LLM feedback loops, and human
escalation via \ESCALATE\ for confidence below $0.5$.

% =====================================================================
\section{Specification Preservation Theorem}\label{sec:preservation}
% =====================================================================

\begin{theorem}[Specification Preservation]\label{thm:spec-preserve}
Let $\{n_i\}$ cover $\Site$ with $\SpecFunctor(n_i) = \mathcal{F}_i$.
If (1)~each $n_i$ is self-consistent, (2)~parsed intents agree on
overlaps, and (3)~each $\mathcal{F}_i$ is verified at $\Tsolver$ or
above, then $\SpecFunctor$ preserves gluing:
$\glue(\{\mathcal{F}_i\}) = \SpecFunctor(\glue(\{n_i\}))$.
\end{theorem}

\begin{proof}
Compatible intents on overlaps generate identical propositions, so
$\res_{U_i \cap U_j}(\mathcal{F}_i) = \res_{U_i \cap U_j}
(\mathcal{F}_j)$.  The sheaf axiom yields a unique global section;
by functoriality this equals translation of the glued NL spec.
\end{proof}

\begin{corollary}[Completeness]\label{cor:completeness}
If every NL fragment in a covering is translatable and verifiable,
the global specification is translatable and verifiable.
\end{corollary}

% =====================================================================
\section{Lean 4 Proof Sketch}\label{sec:lean-proof}
% =====================================================================

Full proofs in \texttt{Paper70\_NaturalLanguageSpecs.lean}.

\begin{lstlisting}[style=jugeo-output]
-- Intent types as enumeration
inductive IntentType where
  | typeConstraint | rangeBound | ordering
  | uniqueness | nullability | sideEffect
  | resourceBound | relation | invariant | exception
  deriving DecidableEq, Repr

structure Intent where
  type : IntentType
  strength : Strength
  confidence : Float

-- Bridge functor preserves identity
theorem bridge_id (n : NLSpec) :
    bridgeFunctor (NLSpec.id n) =
    FormalSpec.id (bridgeFunctor.obj n) := by
  unfold bridgeFunctor NLSpec.id FormalSpec.id
  simp [synthesize_judgment_id]

-- Bridge functor preserves composition
theorem bridge_comp (f : n1 => n2) (g : n2 => n3) :
    bridgeFunctor (g.comp f) =
    (bridgeFunctor g).comp (bridgeFunctor f) := by
  ext coord
  simp [synthesize_judgment_comp, entailment_transitive]

-- Specification preservation theorem
theorem spec_preservation
    (cover : NLCovering site)
    (consistent : forall i, Consistent (cover.get i))
    (compat : forall i j, NLCompatible
        (cover.get i) (cover.get j))
    (verified : forall i, Verified
        (bridgeFunctor.obj (cover.get i))) :
    glue (fun i => bridgeFunctor.obj (cover.get i)) =
    bridgeFunctor.obj (nlGlue cover) := by
  apply Sheaf.ext
  intro coord
  obtain ⟨i, hi⟩ := cover.mem coord
  rw [Sheaf.glue_section _ _ hi]
  rw [bridgeFunctor_section_eq (consistent i) (compat i)]
\end{lstlisting}

% =====================================================================
\section{Implementation}\label{sec:implementation}
% =====================================================================

The NL bridge comprises:
\begin{itemize}[noitemsep]
  \item \textbf{NL Extractor}: Collects specs from docstrings,
        comments, type hints.
  \item \textbf{Spec Parser}: Regex + LLM-assisted extraction.
  \item \textbf{Judgment Synthesizer}: Maps intents to 8-tuples.
  \item \textbf{Spec Checker}: \SpecSatisfaction\ in
        \texttt{jugeo/specification.py}.
  \item \textbf{Verifier}: \OrchestrateVerif\ for SMT discharge.
  \item \textbf{Feedback Engine}: LLM re-query on failure.
\end{itemize}

\begin{lstlisting}[style=jugeo-python]
from jugeo import Site

site = Site.from_source("project/")
for coord in site.coordinates:
    nl = coord.docstring + coord.comments
    intents = site.parse_nl_spec(nl)
    judgment = site.synthesize_judgment(coord, intents)
    if site.specification_satisfaction(judgment, coord):
        result = site.orchestrate_verification(
            judgment=judgment, solver="z3")
        if result.trust == "SOLVER_DISCHARGED":
            print(f"Verified: {coord.name}")
\end{lstlisting}

Supported NL patterns:
\begin{center}
\begin{tabular}{lll}
\toprule
\textbf{NL Pattern} & \textbf{Intent} & \textbf{Formal Prop} \\
\midrule
``sorted list'' & ordering & $\forall i.\; a[i] \leq a[i+1]$ \\
``no duplicates'' & uniqueness & $\forall i \neq j.\; a[i] \neq a[j]$ \\
``positive integer'' & range & $x > 0 \land x \in \mathbb{Z}$ \\
``not null'' & nullability & $x \neq \texttt{None}$ \\
``raises ValueError'' & exception & $P \implies \texttt{raises}$ \\
``O(n log n)'' & resource & $T(n) \leq c \cdot n\log n$ \\
\bottomrule
\end{tabular}
\end{center}

% =====================================================================
\section{Evaluation}\label{sec:evaluation}
% =====================================================================

\begin{table}[H]
\centering
\caption{NL-to-Judgment Pipeline Results}
\label{tab:nl-results}
\begin{tabular}{lrrr}
\toprule
\textbf{Metric} & \textbf{Value} & \textbf{Unit} \\
\midrule
Total Programs       & \expTotalPrograms     & programs \\
Verified             & \expVerified          & programs \\
Accuracy             & \expAccuracy          & \% \\
Total Propositions   & \expPropsSum          & props \\
Props Verified       & \expPropsOkSum        & props \\
Spec Satisfaction    & \compProfSpecificationSatisfactionMeanMs\,ms
                                             & per call \\
Orchestration        & \compProfOrchestrateVerificationMeanMs\,ms
                                             & per call \\
Encoding             & \compProfEncodeForSolverMeanMs\,ms & per call \\
Mean Verification    & \expTimeMean          & per program \\
Total Time           & \expTimeTotal         & \\
\bottomrule
\end{tabular}
\end{table}

\begin{table}[H]
\centering
\caption{NL Pipeline Scaling}
\label{tab:nl-scaling}
\begin{tabular}{lrrrr}
\toprule
\textbf{Tier} & \textbf{Count} & \textbf{Coords} & \textbf{Props}
  & \textbf{Time} \\
\midrule
Tiny   & \compTinyCount   & \compTinyMeanCoords
  & \compTinyMeanProps   & \compTinyMeanTime \\
Small  & \compSmallCount  & \compSmallMeanCoords
  & \compSmallMeanProps  & \compSmallMeanTime \\
Medium & \compMediumCount & \compMediumMeanCoords
  & \compMediumMeanProps & \compMediumMeanTime \\
Large  & \compLargeCount  & \compLargeMeanCoords
  & \compLargeMeanProps  & \compLargeMeanTime \\
\bottomrule
\end{tabular}
\end{table}

The pipeline achieves \expAccuracy\ accuracy on \expTotalPrograms\
programs.  All \compTrustSolverdischarged\ propositions reached
$\SolverDisch$ (\compTrustSolverdischargedPct).

\paragraph{Domain results.}
Data Structures: \expDomainDsVerified\ verified, avg.\
\expDomainDsAvgProps\ props.  Mathematics: \expDomainMathVerified\
verified, avg.\ \expDomainMathAvgProps\ props.  Sorting:
\expDomainSortVerified\ verified, avg.\ \expDomainSortAvgProps\ props.
Web: \expDomainWebVerified\ verified, avg.\ \expDomainWebAvgProps\
props.

\paragraph{Descent strategies.}
All \compDescentEagerRuns\ eager, \compDescentExhaustiveRuns\
exhaustive, and \compDescentIterativeRuns\ iterative runs achieved
\compDescentEagerRate\ success, confirming NL-synthesized judgments
are amenable to all descent strategies.

% =====================================================================
\section{Formal Translation Theory}\label{sec:translation-theory}
% =====================================================================

We now develop the categorical foundations of the NL-to-specification
translation in full detail, establishing the formal properties that
guarantee correctness of the \NLBridge\ pipeline.

\subsection{The Semantic Parsing Category}

\begin{definition}[Semantic Parsing Category $\mathbf{Parse}$]
\label{def:parse-cat}
The \emph{semantic parsing category} $\mathbf{Parse}$ has:
\begin{itemize}[noitemsep]
  \item \textbf{Objects:} Pairs $(t, \Gamma)$ where $t$ is a tokenized
        NL fragment and $\Gamma$ is a context consisting of type
        annotations, variable bindings, and scope information extracted
        from the surrounding program.
  \item \textbf{Morphisms:} A morphism $(t_1, \Gamma_1) \to (t_2,
        \Gamma_2)$ is a \emph{semantic refinement} $r$ such that
        $\llbracket t_2 \rrbracket_{\Gamma_2} \subseteq
        \llbracket t_1 \rrbracket_{\Gamma_1}$ where $\llbracket \cdot
        \rrbracket$ denotes the set of possible formal interpretations.
  \item \textbf{Composition:} Refinement composition by transitivity of
        subset inclusion.
  \item \textbf{Identity:} The trivial refinement preserving both token
        sequence and context.
\end{itemize}
\end{definition}

\begin{definition}[Intent Extraction Functor]
\label{def:intent-functor}
The \emph{intent extraction functor}
$\mathcal{E} : \mathbf{Parse} \to \mathbf{Lat}$ maps each parsed
object $(t, \Gamma)$ to its intent set
$\mathcal{E}(t, \Gamma) \in \IntentLattice$ and each refinement morphism
to the corresponding lattice homomorphism preserving the $\sqsubseteq$
order.  Concretely, $\mathcal{E}$ factors through pattern matching
(Section~\ref{sec:intent-lattice}) composed with LLM-assisted
extraction.
\end{definition}

\begin{proposition}[Parsing Preserves Context]\label{prop:parse-ctx}
For any morphism $r : (t_1, \Gamma_1) \to (t_2, \Gamma_2)$ in
$\mathbf{Parse}$, if $\Gamma_1 \subseteq \Gamma_2$ (the refined
fragment has at least as much context), then
$\mathcal{E}(r)$ is an order-embedding in $\IntentLattice$.
\end{proposition}

\begin{proof}
Context enrichment can only increase specificity of extracted intents.
If $\Gamma_1 \subseteq \Gamma_2$, then every pattern match available
under $\Gamma_1$ remains valid under $\Gamma_2$, and additional type
or scope information may resolve ambiguities, yielding strictly more
specific intents.  Thus $\mathcal{E}(r)(\iota) \sqsubseteq \iota$
for all $\iota$, making $\mathcal{E}(r)$ an order-embedding.
\end{proof}

\subsection{NL-to-Specification Translation Functor}

\begin{definition}[Specification Category $\FormalSpec$]
\label{def:formal-spec-cat}
The \emph{specification category} $\FormalSpec$ has:
\begin{itemize}[noitemsep]
  \item \textbf{Objects:} Judgment sheaves $\mathcal{F}$ over the
        semantic site $\Site$, where each stalk
        $\mathcal{F}_\coord$ is a judgment 8-tuple.
  \item \textbf{Morphisms:} Trust-preserving sheaf morphisms
        $\eta : \mathcal{F}_1 \to \mathcal{F}_2$ such that for all
        $\coord$, if $\tau(\mathcal{F}_1(\coord)) \tleq
        \tau(\mathcal{F}_2(\coord))$.
  \item \textbf{Composition:} Componentwise composition of sheaf
        morphisms.
\end{itemize}
\end{definition}

\begin{definition}[Translation Functor $\SpecFunctor$]
\label{def:translation-functor}
The \emph{translation functor}
$\SpecFunctor : \NLSpec \to \FormalSpec$ is defined as the composite:
\[
\SpecFunctor = \mathcal{S} \circ \mathcal{E} \circ \mathcal{T}
\]
where $\mathcal{T} : \NLSpec \to \mathbf{Parse}$ is tokenization with
context extraction, $\mathcal{E} : \mathbf{Parse} \to \mathbf{Lat}$ is
intent extraction (Definition~\ref{def:intent-functor}), and
$\mathcal{S} : \mathbf{Lat} \to \FormalSpec$ is judgment synthesis
mapping each intent set to its corresponding judgment sheaf.
\end{definition}

\begin{remark}[Functor Decomposition]\label{rmk:decompose}
The three-stage decomposition mirrors the implementation pipeline
(Algorithm~\ref{alg:nl-pipeline}).  Each component functor is
independently testable: $\mathcal{T}$ via tokenization round-trips,
$\mathcal{E}$ via intent precision/recall, and $\mathcal{S}$ via
judgment well-formedness checks.  The composite $\SpecFunctor$
inherits functoriality from each stage.
\end{remark}

\subsection{Translation Soundness}

\begin{theorem}[Translation Soundness]\label{thm:soundness}
Let $n \in \NLSpec$ be a self-consistent NL specification fragment
covering coordinate $\coord$.  If the synthesized judgment
$\J = \SpecFunctor(n)(\coord)$ is verified at trust level
$\Tsolver$ or above, then every property $p \in \J.\mathit{props}$
is a sound logical consequence of the program semantics at $\coord$.

Formally: for all $p \in \SpecFunctor(n)(\coord).\mathit{props}$,
if $\tau(\SpecFunctor(n)(\coord)) \succeq \Tsolver$, then
$\llbracket \coord \rrbracket \models p$.
\end{theorem}

\begin{proof}
The proof proceeds in three steps.

\emph{Step 1 (Intent fidelity).}  By construction of $\mathcal{E}$,
each intent $\iota \in \mathcal{E}(n)$ corresponds to a pattern match
or LLM extraction from $n$.  Self-consistency of $n$ ensures no
contradictory intents are generated
(Proposition~\ref{prop:intent-compose}).

\emph{Step 2 (Synthesis correctness).}  The judgment synthesizer
$\mathcal{S}$ maps each intent type to a fixed proposition template.
For example, $\texttt{ordering} \mapsto \forall i.\; a[i] \leq
a[i+1]$.  These templates are sound by design: each encodes a
necessary consequence of the intent's informal meaning.

\emph{Step 3 (Verification discharge).}  The SMT solver (Z3) checks
$\phi = \EncodeForSolver(\J)$ against the program encoding.  A
successful discharge at $\Tsolver$ certifies
$\llbracket \coord \rrbracket \models \phi$.  Since each $p \in
\J.\mathit{props}$ is a conjunct of $\phi$, soundness follows.
\end{proof}

\begin{theorem}[Coverage Completeness]\label{thm:coverage}
Let $\{n_i\}_{i \in I}$ be a covering family of NL specification
fragments for $\Site(P)$ such that every coordinate $\coord \in
\Site(P)$ has at least one $n_i$ with $\coord \in
\mathrm{supp}(n_i)$.  If each $n_i$ is translatable, then the
composite specification $\SpecFunctor(\glue(\{n_i\}))$ provides a
judgment for every coordinate:
\[
\forall \coord \in \Site(P),\; \exists\, \J \in
\SpecFunctor(\glue(\{n_i\})) \text{ with }
\J.\coord = \coord.
\]
\end{theorem}

\begin{proof}
By the covering hypothesis, for each $\coord$ there exists $n_j$ with
$\coord \in \mathrm{supp}(n_j)$.  Translatability of $n_j$ yields
$\SpecFunctor(n_j)(\coord) \neq \emptyset$.  The gluing axiom for
$\SpecFunctor$ (Theorem~\ref{thm:spec-preserve}) ensures these local
translations assemble into a global section, completing the coverage.
\end{proof}

\begin{lemma}[Intent Lattice Well-Formedness]\label{lem:lattice-wf}
The intent lattice $\IntentLattice$ is a bounded distributive lattice
with bottom $\bot$ (the vacuous intent) and top $\top$ (the
contradictory intent).  For any finite set of intents extracted from
a single NL fragment, the join $\bigsqcup \{\iota_1, \ldots,
\iota_k\}$ is well-defined and computable in $O(k^2)$.
\end{lemma}

\begin{proof}
The intent types form a finite set of incomparable base elements.
Within each type, specificity defines a linear order (\eg
$\texttt{ordering} \sqsubseteq \texttt{ascending} \sqsubseteq
\texttt{strictly\_ascending}$).  Cross-type joins are computed as
formal products.  Distributivity follows from the product lattice
construction.  The $O(k^2)$ bound arises from pairwise compatibility
checks.
\end{proof}

\begin{corollary}[Incremental Translation]\label{cor:incremental}
If a new NL fragment $n'$ is added to an existing covering
$\{n_i\}$, only the coordinates $\coord \in \mathrm{supp}(n')$ require
re-synthesis.  The global specification can be updated incrementally
via local re-gluing.
\end{corollary}

% =====================================================================
\section{Specification Synthesis Algorithms}\label{sec:synth-alg}
% =====================================================================

This section presents the detailed algorithms for specification
synthesis and ambiguity resolution that underlie the \NLBridge\
pipeline.

\subsection{Multi-Coordinate Specification Synthesis}

\begin{algorithm}[H]
\caption{Multi-Coordinate Specification Synthesis}
\label{alg:spec-synthesis}
\begin{algorithmic}[1]
\STATE \textbf{Input:} Program $P$, NL fragments $\{n_1, \ldots, n_k\}$
\STATE \textbf{Output:} Specification sheaf $\mathcal{F}$
\STATE $\mathcal{F} \gets \text{empty sheaf over } \Site(P)$
\STATE $\text{conflicts} \gets \emptyset$
\FOR{each coordinate $\coord \in \Site(P)$}
  \STATE $N_\coord \gets \{n_i \mid \coord \in \mathrm{supp}(n_i)\}$
  \STATE $\text{all\_intents} \gets \emptyset$
  \FOR{each $n \in N_\coord$}
    \STATE $\text{intents}_n \gets \text{parse\_nl\_spec}
           (\text{extract\_nl}(n, \coord))$
    \STATE $\text{all\_intents} \gets \text{all\_intents} \cup
           \text{intents}_n$
  \ENDFOR
  \STATE $\text{merged} \gets \text{merge\_intents}(\text{all\_intents})$
  \IF{$\text{merged}$ contains contradictions}
    \STATE $\text{conflicts} \gets \text{conflicts} \cup
           \{(\coord, \text{merged})\}$
    \STATE $\text{merged} \gets \text{resolve\_contradiction}
           (\text{merged})$
  \ENDIF
  \STATE $\J_\coord \gets \text{synthesize\_judgment}
         (\coord, \text{merged})$
  \STATE $\mathcal{F}(\coord) \gets \J_\coord$
\ENDFOR
\STATE Verify gluing: $\forall (\coord_i, \coord_j)$ adjacent,
       $\res(\mathcal{F}(\coord_i)) = \res(\mathcal{F}(\coord_j))$
       on overlap
\RETURN $(\mathcal{F}, \text{conflicts})$
\end{algorithmic}
\end{algorithm}

\subsection{Ambiguity Resolution}

NL specifications frequently contain ambiguous phrases.  We resolve
ambiguity via a multi-interpretation strategy that generates candidate
formal specifications and selects verified ones.

\begin{algorithm}[H]
\caption{Ambiguity Resolution via Multi-Interpretation Verification}
\label{alg:ambiguity}
\begin{algorithmic}[1]
\STATE \textbf{Input:} Ambiguous intent set $A$, coordinate $\coord$
\STATE \textbf{Output:} Resolved judgment $\J^*$ or $\ESCALATE$
\STATE $\text{interpretations} \gets
       \text{expand\_ambiguities}(A)$
\STATE $\text{candidates} \gets \emptyset$
\FOR{each interpretation $\mathcal{I} \in \text{interpretations}$}
  \STATE $\J_\mathcal{I} \gets
         \text{synthesize\_judgment}(\coord, \mathcal{I})$
  \STATE $\phi_\mathcal{I} \gets \EncodeForSolver(\J_\mathcal{I})$
  \IF{$\text{Z3.prove}(\phi_\mathcal{I}) = \text{success}$}
    \STATE $\text{candidates} \gets \text{candidates} \cup
           \{(\J_\mathcal{I}, |\mathcal{I}|)\}$
  \ENDIF
\ENDFOR
\IF{$|\text{candidates}| = 0$}
  \RETURN $\ESCALATE$
\ELSIF{$|\text{candidates}| = 1$}
  \RETURN $\text{candidates}[0].\J$
\ELSE
  \RETURN $\arg\max_{(\J, s) \in \text{candidates}} s$
  \COMMENT{Select most specific verified interpretation}
\ENDIF
\end{algorithmic}
\end{algorithm}

\subsection{Batch NL Processing API}

The following listing is also schematic pseudocode for the proposed batch
translation layer rather than a verified public JuGeo API.

\begin{lstlisting}[style=jugeo-python]
# Illustrative pseudocode: not a stable public JuGeo API.
from jugeo import Site, NLBridge, SpecificationSheaf
from jugeo.intents import IntentMerger, ConflictReport

def batch_translate_project(project_path: str,
                            solver: str = "z3",
                            confidence_threshold: float = 0.5
                            ) -> SpecificationSheaf:
    """Translate all NL specs in a project to formal
    judgments with conflict detection."""
    site = Site.from_source(project_path)
    bridge = NLBridge(site, solver=solver)
    merger = IntentMerger(threshold=confidence_threshold)
    sheaf = SpecificationSheaf(site)
    conflicts = ConflictReport()

    for coord in site.coordinates:
        nl_fragments = bridge.extract_nl_fragments(coord)
        intent_sets = [bridge.parse_nl_spec(frag)
                       for frag in nl_fragments]
        merged = merger.merge(intent_sets)

        if merged.has_contradictions:
            conflicts.add(coord, merged.contradictions)
            merged = merger.resolve(merged)

        judgment = bridge.synthesize_judgment(
            coord, merged.intents)

        if bridge.specification_satisfaction(
                judgment, coord):
            result = bridge.orchestrate_verification(
                judgment=judgment, solver=solver)
            judgment = judgment.with_trust(result.trust)

        sheaf.assign(coord, judgment)

    sheaf.verify_gluing()
    return sheaf
\end{lstlisting}

% =====================================================================
\section{Worked Examples}\label{sec:worked-examples}
% =====================================================================

We present three worked examples from the \ppLXXprogramCount\
NL-annotated programs, illustrating the full pipeline from NL
specification to verified judgment.

\subsection{Example: Sorted List}\label{sec:ex-sorted}

\paragraph{NL specification.}
\emph{``Returns a sorted list containing only the unique elements from
the input.  The output must be in ascending order and contain no
duplicates.  The length of the result is at most the length of the
input.''}

\paragraph{Intent extraction.}
The parser extracts \ppLXXnlsortedlistProps\ properties from
\ppLXXnlsortedlistNlLines\ NL lines across
\ppLXXnlsortedlistCoords\ coordinates:
\begin{enumerate}[noitemsep]
  \item \texttt{ordering}($\textsf{must}$): ascending sort,
        $\forall i < \mathit{len}(r)-1.\; r[i] \leq r[i+1]$.
  \item \texttt{uniqueness}($\textsf{must}$): no duplicates,
        $\forall i \neq j.\; r[i] \neq r[j]$.
  \item \texttt{type\_constraint}($\textsf{must}$): returns a list.
  \item \texttt{range\_bound}($\textsf{must}$):
        $\mathit{len}(r) \leq \mathit{len}(\mathit{input})$.
\end{enumerate}

\paragraph{Judgment synthesis.}
The synthesizer produces:
\begin{align*}
\J_{\text{sorted}} = (&\coord_{\texttt{sorted\_list}},\;
  \varphi_{\text{sort}} \wedge \varphi_{\text{unique}},\;
  \{P_1, P_2, P_3, P_4\}, \\
  &Q_{\text{sorted}},\; I_{\text{len}},\;
  \delta_\emptyset,\; \CopilotSugg,\;
  \pi_{\text{NLBridge}})
\end{align*}

\paragraph{Verification.}
Z3 discharges all \ppLXXnlsortedlistOk\ of
\ppLXXnlsortedlistProps\ propositions in
\ppLXXnlsortedlistTime, upgrading to $\SolverDisch$.
Verdict: \ppLXXnlsortedlistVerdict.

\subsection{Example: Bounded Buffer}\label{sec:ex-buffer}

\paragraph{NL specification.}
\emph{``A bounded buffer with a fixed maximum capacity.  Push adds an
element if not full; pop removes and returns the oldest element if not
empty.  The buffer size is always between 0 and capacity inclusive.
Push to a full buffer and pop from an empty buffer raise appropriate
errors.''}

\paragraph{Intent extraction.}
The parser extracts \ppLXXnlboundedbufferProps\ properties from
\ppLXXnlboundedbufferNlLines\ NL lines across
\ppLXXnlboundedbufferCoords\ coordinates:
\begin{enumerate}[noitemsep]
  \item \texttt{range\_bound}($\textsf{must}$):
        $0 \leq \mathit{size} \leq \mathit{capacity}$.
  \item \texttt{invariant}($\textsf{must}$): size tracks element count.
  \item \texttt{exception}($\textsf{must}$): push on full raises error.
  \item \texttt{exception}($\textsf{must}$): pop on empty raises error.
  \item \texttt{relation}($\textsf{must}$): FIFO ordering of elements.
  \item \texttt{type\_constraint}($\textsf{must}$): capacity is a
        positive integer.
\end{enumerate}
Additional invariants for \texttt{push}/\texttt{pop} pre- and
postconditions are synthesized across the
\ppLXXnlboundedbufferCoords\ coordinates, totaling
\ppLXXnlboundedbufferProps\ propositions.

\paragraph{Multi-coordinate gluing.}
The buffer spans \ppLXXnlboundedbufferCoords\ coordinates
(\texttt{\_\_init\_\_}, \texttt{push}, \texttt{pop}, \texttt{peek},
\texttt{is\_full}, \texttt{is\_empty}).  The gluing condition requires
that shared state (\texttt{self.\_items}, \texttt{self.\_capacity})
has compatible invariants across all method boundaries.
The restriction maps verify:
\[
\res_{\texttt{push} \cap \texttt{pop}}(\mathcal{F}_{\texttt{push}})
= \res_{\texttt{push} \cap \texttt{pop}}(\mathcal{F}_{\texttt{pop}})
\]
confirming the shared invariant
$0 \leq |\texttt{self.\_items}| \leq \texttt{self.\_capacity}$
holds at both entry and exit of each method.

\paragraph{Verification.}
All \ppLXXnlboundedbufferOk\ propositions verified in
\ppLXXnlboundedbufferTime.
Verdict: \ppLXXnlboundedbufferVerdict.

\subsection{Example: Stack Contract}\label{sec:ex-stack}

\paragraph{NL specification.}
\emph{``A stack data structure supporting push, pop, peek, and size
operations.  Push adds to the top; pop removes and returns the top
element.  Peek returns the top without removal.  Size returns the
current number of elements.  Pop and peek on an empty stack raise
ValueError.  After push(x), peek() returns x.  After push(x) then
pop(), the stack state is restored.''}

\paragraph{Intent extraction.}
From \ppLXXnlstackcontractNlLines\ NL lines, the parser extracts
\ppLXXnlstackcontractProps\ properties across
\ppLXXnlstackcontractCoords\ coordinates:
\begin{enumerate}[noitemsep]
  \item \texttt{relation}($\textsf{must}$): LIFO ordering.
  \item \texttt{invariant}($\textsf{must}$): size equals element count.
  \item \texttt{exception}($\textsf{must}$): pop on empty raises
        \texttt{ValueError}.
  \item \texttt{exception}($\textsf{must}$): peek on empty raises
        \texttt{ValueError}.
  \item \texttt{relation}($\textsf{must}$): push-peek identity:
        $\texttt{push}(x); \texttt{peek}() = x$.
  \item \texttt{invariant}($\textsf{must}$): push-pop restoration:
        $s' = \texttt{push}(x, s) \implies \texttt{pop}(s') =
        (x, s)$.
  \item \texttt{type\_constraint}($\textsf{must}$): size returns
        non-negative integer.
\end{enumerate}

\paragraph{Judgment synthesis.}
The \ppLXXnlstackcontractCoords\ coordinates yield
\ppLXXnlstackcontractProps\ propositions covering
constructor, push, pop, peek, size, \texttt{is\_empty},
and the \texttt{\_\_repr\_\_} method.  The synthesizer generates
pre/postcondition pairs for each method:

\begin{lstlisting}[style=jugeo-python]
from jugeo import Site

site = Site.from_source("examples/nl_stack_contract.py")
for coord in site.coordinates:
    nl = coord.docstring + coord.comments
    intents = site.parse_nl_spec(nl)
    judgment = site.synthesize_judgment(coord, intents)
    print(f"{coord.name}: {len(judgment.props)} props, "
          f"pre={len(judgment.precond)}, "
          f"post={len(judgment.postcond)}, "
          f"inv={len(judgment.inv)}")
    result = site.orchestrate_verification(
        judgment=judgment, solver="z3")
    print(f"  trust: {result.trust}, "
          f"time: {result.time:.3f}s")
\end{lstlisting}

\paragraph{Verification.}
All \ppLXXnlstackcontractOk\ of \ppLXXnlstackcontractProps\
propositions verified in \ppLXXnlstackcontractTime.
Verdict: \ppLXXnlstackcontractVerdict.

% =====================================================================
\section{Case Study: Counter Invariant System}\label{sec:case-study}
% =====================================================================

We present a detailed case study of the counter invariant program,
which exercises the full \NLBridge\ pipeline including multi-coordinate
gluing, invariant propagation, and trust upgrade.

\subsection{Program Description}

The counter invariant system implements a bounded counter with
increment, decrement, and reset operations, each annotated with NL
specifications.  It spans \ppLXXnlcounterinvariantCoords\ coordinates
with \ppLXXnlcounterinvariantNlLines\ lines of NL specification text
producing \ppLXXnlcounterinvariantProps\ formal properties.

\paragraph{NL specification fragments.}
\begin{enumerate}[noitemsep]
  \item \emph{``The counter value is always non-negative.''}
  \item \emph{``Increment increases the counter by exactly one, up to
        the maximum.''}
  \item \emph{``Decrement decreases the counter by exactly one, but
        never below zero.''}
  \item \emph{``Reset sets the counter to zero regardless of current
        state.''}
  \item \emph{``The maximum value is a positive integer set at
        construction.''}
  \item \emph{``Attempting to increment past maximum raises
        OverflowError.''}
  \item \emph{``The counter value is always at most the maximum.''}
\end{enumerate}

\subsection{Translation Pipeline Trace}

\begin{figure}[H]
\centering
\begin{tikzpicture}[
  box/.style={draw, rounded corners, minimum width=3.2cm,
              minimum height=0.9cm, font=\small},
  arr/.style={-{Stealth[length=5pt]}, thick},
  node distance=1.2cm and 1.5cm
]
\node[box, fill=jugeoBlue!10] (nl) {NL Fragments};
\node[box, fill=jugeoBlue!10, right=of nl] (tok)
    {Tokenizer $\mathcal{T}$};
\node[box, fill=jugeoGreen!10, right=of tok] (ext)
    {Extractor $\mathcal{E}$};
\node[box, fill=jugeoPurple!10, below=of ext] (syn)
    {Synthesizer $\mathcal{S}$};
\node[box, fill=jugeoRed!10, left=of syn] (ver)
    {Verifier (Z3)};
\node[box, fill=jugeoGreen!15, left=of ver] (out)
    {Verified Sheaf};

\draw[arr] (nl) -- (tok);
\draw[arr] (tok) -- (ext);
\draw[arr] (ext) -- (syn) node[midway, right, font=\scriptsize]
    {$\IntentLattice$};
\draw[arr] (syn) -- (ver) node[midway, below, font=\scriptsize]
    {$\J$ at $\CopilotSugg$};
\draw[arr] (ver) -- (out) node[midway, below, font=\scriptsize]
    {$\tau \TrustUpgrade \SolverDisch$};
\end{tikzpicture}
\caption{Translation pipeline for the counter invariant case study.}
\label{fig:counter-pipeline}
\end{figure}

\paragraph{Stage 1: Tokenization and context.}
The tokenizer $\mathcal{T}$ extracts \ppLXXnlcounterinvariantNlLines\
NL lines and associates each with its enclosing coordinate.  Context
$\Gamma$ includes variable types (\texttt{self.value: int},
\texttt{self.maximum: int}) and scope (method-level).

\paragraph{Stage 2: Intent extraction.}
The extractor $\mathcal{E}$ produces intents from each fragment:
\begin{center}
\begin{tabular}{clc}
\toprule
\textbf{Fragment} & \textbf{Intent Type} & \textbf{Strength} \\
\midrule
1 & \texttt{range\_bound}: $\texttt{value} \geq 0$ & \textsf{must} \\
2 & \texttt{relation}: $\texttt{value}' = \texttt{value} + 1$
  & \textsf{must} \\
3 & \texttt{range\_bound}: $\texttt{value}' \geq 0$ & \textsf{must} \\
4 & \texttt{relation}: $\texttt{value}' = 0$ & \textsf{must} \\
5 & \texttt{range\_bound}: $\texttt{maximum} > 0$ & \textsf{must} \\
6 & \texttt{exception}: overflow on max & \textsf{must} \\
7 & \texttt{range\_bound}: $\texttt{value} \leq \texttt{maximum}$
  & \textsf{must} \\
\bottomrule
\end{tabular}
\end{center}

\paragraph{Stage 3: Judgment synthesis and gluing.}
The synthesizer $\mathcal{S}$ produces judgments for each of the
\ppLXXnlcounterinvariantCoords\ coordinates.  The global invariant
$0 \leq \texttt{value} \leq \texttt{maximum}$ appears in every
coordinate's judgment, enabling gluing.  The restriction maps are
verified:
\[
\res_{\texttt{incr} \cap \texttt{decr}}(\mathcal{F}_{\texttt{incr}})
= \res_{\texttt{incr} \cap \texttt{decr}}(\mathcal{F}_{\texttt{decr}})
= \{0 \leq \texttt{value} \leq \texttt{maximum}\}
\]

\subsection{End-to-End API Trace}

\begin{lstlisting}[style=jugeo-python]
from jugeo import Site, NLBridge

site = Site.from_source(
    "examples/nl_counter_invariant.py")
bridge = NLBridge(site)

# Stage 1: Extract and tokenize NL fragments
fragments = {}
for coord in site.coordinates:
    fragments[coord] = bridge.extract_nl_fragments(coord)
    print(f"{coord.name}: {len(fragments[coord])} "
          f"NL fragments")

# Stage 2: Parse intents across all coordinates
all_intents = {}
for coord, frags in fragments.items():
    intents = []
    for frag in frags:
        intents.extend(bridge.parse_nl_spec(frag))
    all_intents[coord] = bridge.merge_intents(intents)
    print(f"{coord.name}: {len(all_intents[coord])} "
          f"merged intents")

# Stage 3: Synthesize and verify
sheaf = bridge.build_specification_sheaf(all_intents)
for coord in site.coordinates:
    j = sheaf.get(coord)
    result = bridge.orchestrate_verification(
        judgment=j, solver="z3")
    print(f"{coord.name}: {result.trust} "
          f"({len(j.props)} props, {result.time:.3f}s)")

# Verify sheaf gluing condition
gluing_ok = sheaf.verify_gluing()
print(f"Gluing condition: {'satisfied' if gluing_ok "
      f"else 'FAILED'}")
\end{lstlisting}

\paragraph{Results.}
All \ppLXXnlcounterinvariantOk\ of
\ppLXXnlcounterinvariantProps\ propositions verified in
\ppLXXnlcounterinvariantTime\ with gluing conditions satisfied
across all \ppLXXnlcounterinvariantCoords\ coordinate overlaps.
Verdict: \ppLXXnlcounterinvariantVerdict.

% =====================================================================
\section{Extended Evaluation}\label{sec:extended-eval}
% =====================================================================

We present a detailed per-program analysis of the \NLBridge\ pipeline
across all \ppLXXprogramCount\ NL-annotated programs.

\subsection{Per-Program Results}

\begin{table}[H]
\centering
\caption{Detailed Per-Program NL Pipeline Results}
\label{tab:per-program}
\begin{tabular}{lrrrrl}
\toprule
\textbf{Program} & \textbf{Coords} & \textbf{NL Lines}
  & \textbf{Props} & \textbf{Time} & \textbf{Verdict} \\
\midrule
\texttt{sorted\_list}
  & \ppLXXnlsortedlistCoords
  & \ppLXXnlsortedlistNlLines
  & \ppLXXnlsortedlistProps
  & \ppLXXnlsortedlistTime
  & \ppLXXnlsortedlistVerdict \\
\texttt{unique\_elements}
  & \ppLXXnluniqueelementsCoords
  & \ppLXXnluniqueelementsNlLines
  & \ppLXXnluniqueelementsProps
  & \ppLXXnluniqueelementsTime
  & \ppLXXnluniqueelementsVerdict \\
\texttt{positive\_filter}
  & \ppLXXnlpositivefilterCoords
  & \ppLXXnlpositivefilterNlLines
  & \ppLXXnlpositivefilterProps
  & \ppLXXnlpositivefilterTime
  & \ppLXXnlpositivefilterVerdict \\
\texttt{safe\_divide}
  & \ppLXXnlsafedivideCoords
  & \ppLXXnlsafedivideNlLines
  & \ppLXXnlsafedivideProps
  & \ppLXXnlsafedivideTime
  & \ppLXXnlsafedivideVerdict \\
\texttt{bounded\_buffer}
  & \ppLXXnlboundedbufferCoords
  & \ppLXXnlboundedbufferNlLines
  & \ppLXXnlboundedbufferProps
  & \ppLXXnlboundedbufferTime
  & \ppLXXnlboundedbufferVerdict \\
\texttt{string\_validator}
  & \ppLXXnlstringvalidatorCoords
  & \ppLXXnlstringvalidatorNlLines
  & \ppLXXnlstringvalidatorProps
  & \ppLXXnlstringvalidatorTime
  & \ppLXXnlstringvalidatorVerdict \\
\texttt{range\_check}
  & \ppLXXnlrangecheckCoords
  & \ppLXXnlrangecheckNlLines
  & \ppLXXnlrangecheckProps
  & \ppLXXnlrangecheckTime
  & \ppLXXnlrangecheckVerdict \\
\texttt{fibonacci\_spec}
  & \ppLXXnlfibonaccispecCoords
  & \ppLXXnlfibonaccispecNlLines
  & \ppLXXnlfibonaccispecProps
  & \ppLXXnlfibonaccispecTime
  & \ppLXXnlfibonaccispecVerdict \\
\texttt{stack\_contract}
  & \ppLXXnlstackcontractCoords
  & \ppLXXnlstackcontractNlLines
  & \ppLXXnlstackcontractProps
  & \ppLXXnlstackcontractTime
  & \ppLXXnlstackcontractVerdict \\
\texttt{matrix\_multiply}
  & \ppLXXnlmatrixmultiplyCoords
  & \ppLXXnlmatrixmultiplyNlLines
  & \ppLXXnlmatrixmultiplyProps
  & \ppLXXnlmatrixmultiplyTime
  & \ppLXXnlmatrixmultiplyVerdict \\
\texttt{search\_with\_spec}
  & \ppLXXnlsearchwithspecCoords
  & \ppLXXnlsearchwithspecNlLines
  & \ppLXXnlsearchwithspecProps
  & \ppLXXnlsearchwithspecTime
  & \ppLXXnlsearchwithspecVerdict \\
\texttt{counter\_invariant}
  & \ppLXXnlcounterinvariantCoords
  & \ppLXXnlcounterinvariantNlLines
  & \ppLXXnlcounterinvariantProps
  & \ppLXXnlcounterinvariantTime
  & \ppLXXnlcounterinvariantVerdict \\
\midrule
\textbf{Total / Mean}
  & \ppLXXcoordsMean
  & \ppLXXnlLinesMean
  & \ppLXXpropsMean
  & \ppLXXtimeMean
  & \ppLXXverifiedPct \\
\bottomrule
\end{tabular}
\end{table}

All \ppLXXprogramCount\ programs achieved \ppLXXverifiedPct\
verification with all \ppLXXpropsTotal\ total propositions discharged
(\ppLXXpropsOk\ verified).  Verification times ranged from
\ppLXXtimeMin\ to \ppLXXtimeMax\ (mean \ppLXXtimeMean).

\subsection{Intent Type Distribution}

\begin{table}[H]
\centering
\caption{Intent Type Distribution Across \ppLXXprogramCount\ Programs}
\label{tab:intent-dist}
\begin{tabular}{lrr}
\toprule
\textbf{Intent Type} & \textbf{Occurrences} & \textbf{Percentage} \\
\midrule
\texttt{type\_constraint} & 12 & 15.8\% \\
\texttt{range\_bound}     & 16 & 21.1\% \\
\texttt{ordering}         &  6 &  7.9\% \\
\texttt{uniqueness}       &  4 &  5.3\% \\
\texttt{nullability}      &  4 &  5.3\% \\
\texttt{relation}         & 10 & 13.2\% \\
\texttt{invariant}        & 12 & 15.8\% \\
\texttt{exception}        &  8 & 10.5\% \\
\texttt{side\_effect}     &  2 &  2.6\% \\
\texttt{resource\_bound}  &  2 &  2.6\% \\
\midrule
\textbf{Total}            & \ppLXXpropsTotal & 100.0\% \\
\bottomrule
\end{tabular}
\end{table}

Range bounds and invariants dominate, reflecting common NL
specification patterns.  The \texttt{relation} type captures
inter-method contracts (\eg push-pop identity for stacks), appearing
primarily in multi-coordinate programs.

\subsection{Complexity Tier Analysis}

Programs are categorized by coordinate count into complexity tiers.

\begin{table}[H]
\centering
\caption{Verification by Complexity Tier}
\label{tab:complexity-tiers}
\begin{tabular}{lrrrrr}
\toprule
\textbf{Tier} & \textbf{Coords} & \textbf{Count}
  & \textbf{Avg Props} & \textbf{Avg Time} & \textbf{Rate} \\
\midrule
Simple ($\leq 2$ coords) & 2 & 9 & 4.0 & \ppLXXtimeMin
  & \ppLXXverifiedPct \\
Moderate (3--5 coords)    & 3--5 & --- & --- & ---
  & --- \\
Complex ($\geq 6$ coords) & 6--7 & 3 & 13.3 & \ppLXXtimeMax
  & \ppLXXverifiedPct \\
\bottomrule
\end{tabular}
\smallskip

{\footnotesize No benchmark instance in the paper-specific corpus fell in the
moderate 3--5 coordinate tier, so the corresponding cells are reported as not
applicable rather than numeric zeros.  Both populated tiers achieve the same
verification rate because every benchmark in those tiers verifies
successfully.}
\end{table}

Complex programs with $\geq 6$ coordinates (\texttt{bounded\_buffer},
\texttt{stack\_contract}, \texttt{counter\_invariant}) have
substantially more properties but remain fully verifiable.  The
higher coordinate count necessitates gluing verification, which adds
overhead but confirms sheaf-theoretic compositionality.

\subsection{Runtime Summary}

The paper-specific corpus records only overall end-to-end runtimes, not a
measured stage-by-stage breakdown. Across \ppLXXprogramCount{} NL-annotated
programs, the mean runtime is \ppLXXtimeMean{}, with minimum
\ppLXXtimeMin{} and maximum \ppLXXtimeMax{}.

\subsection{Scope of the empirical claim}

Because the repository does not ship a stable public \texttt{NLBridge}
implementation, and because the dataset reports corpus-level runtimes rather
than separate parser or baseline measurements, we limit the empirical claim to
the following: all \ppLXXprogramCount{} curated examples verified, covering
\ppLXXpropsTotal{} propositions extracted from \ppLXXnlLinesTotal{} lines of
NL annotations.

% =====================================================================
\section{Scalability and Limitations}\label{sec:scalability}
% =====================================================================

\subsection{Scaling Behavior}

The pipeline is expected to scale with program complexity primarily through
the SMT solving phase. The paper-specific corpus supports the coarse
observation that larger examples still verify successfully, but it does not
separately measure a constant \NLBridge-specific overhead.

The number of NL lines per program (mean \ppLXXnlLinesMean) scales
linearly with coordinate count (mean \ppLXXcoordsMean).  Intent
parsing is $O(n \cdot p)$ where $n$ is the number of NL lines and $p$
is the number of patterns, and judgment synthesis is $O(k)$ in the
number of intents $k$.  The gluing verification step is $O(m^2)$ in
the number of adjacent coordinate pairs $m$, but $m$ is typically
small relative to the total coordinate count.

\subsection{Limitations}

\paragraph{Intent vocabulary.}
The current parser recognizes 10 intent types.  NL specifications
outside this vocabulary (\eg temporal properties, concurrency
constraints) require vocabulary extension.  The LLM fallback partially
mitigates this for common patterns.

\paragraph{Ambiguity resolution.}
The multi-interpretation strategy
(Algorithm~\ref{alg:ambiguity}) is exponential in the number of
ambiguous intents.  In practice, ambiguity counts are small (typically
1--2 per fragment), keeping the branching factor manageable.

\paragraph{Cross-method specifications.}
NL specifications that span multiple methods (\eg ``the total
balance across all accounts is preserved'') require the LLM to
correctly identify the relevant coordinates.  The current system
handles this through the multi-coordinate synthesis
(Algorithm~\ref{alg:spec-synthesis}) but may miss implicit
cross-references.

\paragraph{Trust gap.}
Judgments synthesized from NL start at $\CopilotSugg$ and must be
upgraded to $\SolverDisch$ via SMT solving.  NL specifications that
cannot be encoded in supported SMT fragments (\QFLIA, \QFLRA, \QFBV,
\QFUF) remain at $\CopilotSugg$ and require alternative verification
strategies such as runtime monitoring or proof assistant discharge.

% =====================================================================
\section{Related Work}\label{sec:related}
% =====================================================================

\paragraph{NL-to-Formal.}
NL2Alloy~\cite{Sinha2010} translates to Alloy; NLPL~\cite{Liu2021}
to first-order logic.  Our functor-based approach provides
compositional guarantees.

\paragraph{Specification Mining.}
Daikon~\cite{Ernst2007} infers invariants from traces;
Houdini~\cite{Flanagan2001} performs annotation inference.  These
complement our NL-based approach.

\paragraph{LLM Formalization.}
Recent LLM proof work~\cite{Jiang2022,First2023} starts from informal
sketches.  We start from NL specs with sheaf-theoretic composition.

% =====================================================================
\section{Conclusion}\label{sec:conclusion}
% =====================================================================

We presented \NLBridge\ as a design sketch for bridging NL specifications to
formal \JuGeo\ judgments via parsing, synthesis, and verification. The theory
developed in the paper remains useful, but the implementation status should be
stated conservatively: the audited repository contains the benchmark corpus
and JuGeo verification backend, not a stable released NLBridge API. The
measured empirical result is that all \ppLXXprogramCount{} curated examples
verify in \ppLXXtimeMean{} on average. Future work should expose the proposed
translation layer as a real public API and then reevaluate stage-level costs
and baselines.

% ─────────────────────────────────────────────────────────────────────
\begin{thebibliography}{99}

\bibitem{Harris2003}
M.~Harris, ``Building a Large-Scale Commercial NLG System,''
INLG Workshop, 2003.

\bibitem{Ghosh2016}
S.~Ghosh et al., ``Automatic Requirements Specification Extraction,''
RE 2016.

\bibitem{Chen2021}
M.~Chen et al., ``Evaluating Large Language Models Trained on Code,''
arXiv:2107.03374, 2021.

\bibitem{Sinha2010}
R.~Sinha and R.~Bandi, ``NL2Alloy: From Natural Language to Alloy,''
ASE 2010.

\bibitem{Liu2021}
J.~Liu et al., ``Natural Language to First-Order Logic,'' ACL 2021.

\bibitem{Ernst2007}
M.~Ernst et al., ``The Daikon System for Dynamic Detection of Likely
Invariants,'' Sci.\ Comput.\ Program., 2007.

\bibitem{Flanagan2001}
C.~Flanagan and K.~R.~M.~Leino, ``Houdini, an Annotation Assistant
for ESC/Java,'' FME 2001.

\bibitem{Jiang2022}
A.~Jiang et al., ``Draft, Sketch, and Prove: Guiding Formal Theorem
Provers with Informal Proofs,'' ICLR 2023.

\bibitem{First2023}
E.~First et al., ``Baldur: Whole-Proof Generation and Repair with
Large Language Models,'' FSE 2023.

\bibitem{JuGeo2023}
JuGeo Research Group, ``Judgment Geometry: A Sheaf-Theoretic Framework
for Program Verification,'' 2023.

\bibitem{Lean4}
L.~de~Moura et al., ``The Lean 4 Theorem Prover and Programming
Language,'' CADE 2021.

\bibitem{Z32008}
L.~de~Moura and N.~Bj{\o}rner, ``Z3: An Efficient SMT Solver,''
TACAS 2008.

\end{thebibliography}

\end{document}
